% STRUTTURA DEL DOCUMENTO INTERO
\begin{document}
    \frontmatter
    \begin{titlepage}
    \begin{center}
        \begin{LARGE}
            \textbf{\myUni}\\
        \end{LARGE}

        \vspace{10pt}

        \begin{Large}
            \textsc{\myDepartment}\\
        \end{Large}

        \vspace{10pt}

        \begin{large}
            \textsc{\myFaculty}\\
        \end{large}

        % \vspace{30pt}
        \vfill

        \begin{figure}[htbp]
            \centering
            
\includegraphics[width=0.5\textwidth]{unipd-logo} %va bene così perchè le impostazioni di compilazione prevedono che le immagini siano nella cartella "images"
        \end{figure}
        
        %\vspace{30pt}
        \vfill

        \begin{LARGE}
            \textbf{\myTitle}\\
        \end{LARGE}

        \vspace{10pt}

        \begin{large}
            \textsl{\myDegree}\\
        \end{large}

        %\vspace{40pt}
        \vfill

        \begin{large}
            \begin{flushleft}
                \textit{Relatore}\\
                \vspace{5pt}
                \profTitle\ \myProf
            \end{flushleft}

            % You can tweak the spacing to have professor and student names on the same line
            % useful if the page is broken by a long thesis title and you need more space
            % \vspace{-52pt}

            \begin{flushright}
                \textit{Laureanda}\\
                \vspace{5pt}
                \myName \\
                \vspace{5pt}
                \textit{Matricola} \myID
            \end{flushright}
        \end{large}

        %\vspace{40pt}
        \vfill

        \line(1, 0){338} \\
        \begin{normalsize}
            \textsc{Anno Accademico \myAA}
        \end{normalsize}
    \end{center}
\end{titlepage}
 % pagina iniziale con titolo e logo
    \input{preface/copyright} % riga del copywright

    \cleardoublepage % forza a iniziare la prossima sezione a pagina dispari - SOLO IN MODALITà FRONTE RETRO x la stampa
    \cleardoublepage
\phantomsection
\thispagestyle{empty}
\pdfbookmark{Dedica}{Dedica}

\vspace*{5cm}

\begin{flushright}
    \textit{A chi non si arrende,\\
    a chi combatte in silenzio.\\
    A chi ha trovato la forza di rialzarsi da solo,\\
    a chi si è lasciato alle spalle chi non era più.\\
    A chi ha imparato a essere il proprio punto di riferimento:\\
    continua, credici, resisti.\\
    Il traguardo esiste, ed è più vicino di quanto pensi.}
\end{flushright}

%\medskip
%\begin{center}
%    Lorem ipsum dolor sit amet, consectetuer adipiscing elit. \\ \medskip
%    --- Citazione
%\end{center} % dedica
    \cleardoublepage
    \cleardoublepage
\phantomsection
\pdfbookmark{Sommario}{Sommario}
\begingroup
\let\clearpage\relax
\let\cleardoublepage\relax
\let\cleardoublepage\relax

\chapter*{Sommario}
Il presente elaborato documenta lo stage di 320 ore svolto dalla studentessa Giulia Barzon presso Wintech S.p.A., a conclusione del Corso di Laurea in Informatica dell'Università degli Studi di Padova.

L'obiettivo principale dello stage è stato l'analisi, la valutazione e l'integrazione di una piattaforma di \textit{Data Catalog} all'interno dell'infrastruttura aziendale, 
per mantenere un catalogo dati centralizzato e garantire il rispetto delle normative sulla privacy, come il GDPR, e delle policy aziendali.\\

Partendo dall'analisi dei requisiti aziendali e da uno studio comparativo delle principali piattaforme di \textit{Data Catalog} 
presenti sul mercato, è stata individuata la piattaforma open-source OpenMetadata come lo strumento più idoneo a soddisfare le esigenze di Wintech.

La fase successiva si è concentrata sul deployment e sulla configurazione del software, testandone in modo completo le funzionalità chiave: la \textit{Data Ingestion} 
dalle diverse fonti aziendali, la \textit{Data Discovery} e la \textit{Data Quality}. 
Particolare attenzione è stata rivolta alla sicurezza e alla \textit{Data Governance}, validando l'efficacia del motore di individuazione e classificazione automatica dei dati sensibili e personali.\\

Per completare il lavoro svolto, parte delle ore dello stage è stata dedicata allo studio di Microsoft Power BI e dei principi della \textit{Business Intelligence}, 
dimostrando come l'integrazione tra strumenti di BI e piattaforme di \textit{Data Catalog} possa migliorare la gestione dei dati aziendali, garantendo la loro qualità, sicurezza e tracciabilità.

%\vfill

%\selectlanguage{english}
%\pdfbookmark{Abstract}{Abstract}
%\chapter*{Abstract}

%\selectlanguage{italian}

\endgroup

\vfill
 % sommario
    \cleardoublepage
    \cleardoublepage
\pdfbookmark{\contentsname}{tableofcontents}
\setcounter{tocdepth}{2}
\tableofcontents
%\markboth{\contentsname}{\contentsname}

\cleardoublepage
% Figures list
\phantomsection
\pdfbookmark{\listfigurename}{lof}
\listoffigures

\cleardoublepage
% Tables list
\phantomsection
\pdfbookmark{\listtablename}{lot}
\listoftables

\cleardoublepage
 % indice
    \cleardoublepage

    % contenuto principale: capitoli
    \mainmatter
    % capitolo 1: introduzione
\chapter{Contesto aziendale}
\label{cap:contesto-aziendale}

\intro{Il presente capitolo ha lo scopo di introdurre il contesto aziendale in cui si è svolta l'esperienza di stage. Verrà presentata la realtà di Wintech S.p.A. esplorandone 
l'offerta, l'approccio alla sicurezza informatica e il modello organizzativo.}\\

\section{L'azienda}
Wintech (Winning Technologies S.p.A.) è una consolidata realtà italiana operante nel settore dell'Information Technology.\footcite{wintech:website}
Fondata nel 1987, l'azienda ha la sua sede principale a Padova e presidia strategicamente il territorio con ulteriori filiali operative a Milano e Bassano del Grappa.

Ad oggi l'azienda conta un organico di oltre 95 professionisti, un portafoglio di più di 500 clienti attivi e un fatturato superiore ai 12 milioni di euro.
Wintech si posiziona sul mercato principalmente come consulente e \emph{System Integrator}: il suo core business consiste nel combinare competenze e tecnologie avanzate per far comunicare in 
sinergia molteplici sottosistemi eterogenei, creando un'infrastruttura IT coesa, sicura e funzionale che porti un reale valore aggiunto al business del cliente.

I prodotti e i servizi offerti si rivolgono a un bacino d'utenza ampio e diversificato, studiato per rispondere alle esigenze di liberi professionisti, Piccole e Medie Imprese (PMI), 
grandi aziende, Pubblica Amministrazione, istituti bancari e assicurativi. 

Per garantire lo sviluppo e l'integrazione di soluzioni all'avanguardia, stabili e scalabili, Wintech ha stretto nel tempo solide \emph{partnership} con i principali colossi mondiali 
dell'informatica, tra cui Microsoft, IBM, Oracle, HP e Sophos. Tali collaborazioni strategiche garantiscono l'accesso diretto alle migliori tecnologie sul mercato e certificano 
l'elevato livello di competenza raggiunto dal personale aziendale.

Infine, un aspetto cardine dell'identità aziendale è la profonda attenzione rivolta alla sicurezza dei dati e alla qualità dei processi interni. 
A testimonianza di questo impegno, Wintech è ufficialmente dotata delle prestigiose certificazioni internazionali ISO 9001, inerente ai sistemi di gestione per la qualità, 
e ISO/IEC 27001, riguardante i sistemi di gestione per la sicurezza delle informazioni.


% logo Wintech



\section{L'offerta}
L'offerta di Wintech si articola in due macro-aree complementari: i servizi di consulenza e progettazione IT, e la fornitura di applicazioni gestionali specifiche per i diversi settori di mercato.

\subsection{Servizi}
L'azienda accompagna i clienti in ogni fase del loro percorso di transizione digitale. I servizi offerti spaziano dall'analisi e progettazione di soluzioni su misura, 
al \emph{deployment} e all'avviamento dei sistemi, fino alla formazione del personale. Viene inoltre garantita una costante assistenza tecnica e applicativa, sia tramite \emph{help desk} 
remoto che con interventi diretti sul posto. Un pilastro fondamentale dei servizi Wintech è la consulenza in ambito Cyber Security e architetture Cloud, per garantire la protezione, 
il monitoraggio e il Disaster Recovery delle infrastrutture di rete.

\subsection{Prodotti e applicazioni gestionali}
Wintech personalizza e fornisce un ampio portafoglio di \emph{software}, adattandoli alle dimensioni e alle esigenze operative del cliente. L'offerta si suddivide principalmente in:
\begin{itemize}
    \item \textbf{Professionisti e studi}: Soluzioni integrate per l'amministrazione, la fatturazione, la contabilità e la gestione avanzata del personale (tra cui le suite PROFIS, JOB e STUDIO).
    \item \textbf{PMI e grandi aziende}: Sistemi avanzati di Enterprise Resource Planning (ERP) per il controllo di gestione (come eSOLVER), strumenti di Business Intelligence per l'analisi dei dati, e piattaforme per il Customer Relationship Management (CRM).
    \item \textbf{Aziende di moda}: Declinazioni specifiche dei gestionali (es. eSOLVER per la Moda) sviluppate per gestire le varianti di prodotto, la logistica, le spedizioni e i punti vendita tipici del settore.
    \item \textbf{Banche e assicurazioni}: Soluzioni dedicate alla sicurezza documentale, all'archiviazione e all'integrazione della firma grafometrica.
\end{itemize}

A completamento dell'offerta, l'azienda progetta e sviluppa portali web e piattaforme di \emph{Collaboration} aziendale (intranet) per favorire la comunicazione interna e la condivisione delle informazioni.

\subsection{Partner tecnologici}
Per garantire standard di qualità eccellenti e fornire tecnologie sempre all'avanguardia, Wintech si avvale della collaborazione dei principali colossi mondiali del settore IT. 
Le solide \emph{partnership} instaurate nel tempo permettono all'azienda di accedere in modo diretto a soluzioni certificate, scalabili e sicure.

\begin{figure}[H]
    \centering
    % loghi dei partner principali
    
\includegraphics[width=1\textwidth]{partner.png}
    \vspace{0.2cm}
    \caption{Principali \emph{partner} di Wintech}
    \label{fig:aziende_partner}
\end{figure}

\section{Certificazioni aziendali}
Wintech è un'azienda certificata ISO 27001 e ISO 9001, che garantiscono massimi standard di qualità professionale e di gestione della sicurezza ai collaboratori, nonchè reali vantaggi 
competitivi ai clienti.
Tutte le certificazioni:
ISO 9001
ISO IEC 27001
ISO IEC 27017
ISO IEC 27018
ISO 22301
% da descrivere in molto breve

% ISO 9001 (Qualità) e ISO/IEC 27001 (Sicurezza delle informazioni)
% immagine/infografica sulla sicurezza dei dati (da prendere sul sito)


\section{Organizzazione e processi interni}
% organizzazione aziendale: team ()

% processi e way of working interno
% metodologia agile (framework scrum) - con immagine, microsoft teams, regole sulla comunicazione
% Un aspetto positivo dello stage è stata la comunicazione con i colleghi e con la tutor interna, che mi ha fornito supporto quotidiano e con cui ci sono stati confronti aperti



\section{Organizzazione del testo}
Il presente documento è strutturato in nove capitoli, organizzati come segue:

\begin{itemize}
    \item \textbf{Il primo capitolo} introduce il contesto aziendale di Wintech S.p.A..
    \item \textbf{Il secondo capitolo} fornisce una descrizione dettagliata dell'esperienza di stage.
    \item \textbf{Il terzo capitolo} si focalizza sullo studio dei casi d'uso e l'analisi dei requisiti aziendali.
    \item \textbf{Il quarto capitolo} analizza e confronta la piattaforma OpenMetadata con i principali competitor sul mercato in ambito Data Catalog.
    \item \textbf{Il quinto capitolo} è dedicato all'implementazione tecnica, illustrando le fasi di installazione, configurazione e Data Discovery.
    \item \textbf{Il sesto capitolo} approfondisce i concetti di Data Quality e Data Observability per il monitoraggio dello stato di salute dei dati.
    \item \textbf{Il settimo capitolo} tratta il tema della Data Governance e dello sviluppo sicuro.
    \item \textbf{L'ottavo capitolo} approfondisce l'integrazione tra il Data Catalog e gli strumenti di Business Intelligence, in particolare Microsoft Power BI.
    \item \textbf{Il nono capitolo} presenta le conclusioni, con un consuntivo finale sul raggiungimento degli obiettivi, le problematiche affrontate e le competenze acquisite.
\end{itemize}

\subsection{Convenzioni}
Riguardo la stesura del testo, relativamente al documento sono state adottate le seguenti convenzioni tipografiche:
\begin{itemize}
	\item gli acronimi, le abbreviazioni e i termini ambigui o di uso non comune menzionati vengono definiti nel glossario, situato alla fine del presente documento.
	\item per la prima occorrenza dei termini riportati nel glossario viene utilizzata la seguente nomenclatura: \emph{parola}\glsfirstoccur.
	\item particolari termini in lingua straniera o facenti parti del gergo tecnico sono evidenziati con il carattere \emph{corsivo}. Fanno eccezione i macro-concetti o i domini disciplinari del progetto (come Data Quality, Data Catalog, Data Lineage...), 
    i quali vengono scritti in tondo con le iniziali maiuscole per agevolare la scorrevolezza della lettura.
\end{itemize}
 %1
    % capitolo 2: descrizione generale dello stage
\chapter{Progetto di stage}
\label{cap:progetto-di-stage}

\intro{In questo capitolo viene introdotto il progetto di stage, descrivendone le motivazioni iniziali, gli obiettivi prefissati, i vincoli e la pianificazione delle fasi di lavoro.}\\

\section{L'idea}

\subsection{Motivazioni: perchè adottare un Data Catalog}
Nell'era moderna, le aziende hanno un'enormità di dati sparsi in sistemi diversi (database, fogli Excel, file CSV...). Senza un catalogo, le persone 
non sanno cosa esiste, dove si trova o se possono fidarsi di quel dato. Il Data Catalog risolve questo caos creando una mappa unica, organizzata e facilmente consultabile.

\subsubsection{Sicurezza e compliance normativa}
Storicamente l'individuazione dei dati personali ai fini GDPR avveniva in modo manuale, producendo un semplice "catalogo grezzo". 
L'adozione di un Data Catalog moderno automatizza questo processo scansionando i database e classificando i dati sensibili in autonomia.
Il passaggio da un approccio manuale a un catalogo completo e automatizzato segna una forte crescita della maturità aziendale nella sicurezza e nel modello \gls{samm}. 
Grazie a questa mappatura esatta e dinamica, l'azienda sa sempre dove e come applicare le necessarie policy di protezione, come il mascheramento o la pseudonimizzazione.

\subsubsection{Analisi dei dati e Business Intelligence}
Gli analisti accedono al catalogo per cercare le informazioni di cui hanno bisogno per creare i loro report. 
Se un analista ha bisogno del fatturato del 2023, cerca "Fatturato" nel Data Catalog: il sistema gli dice subito in quale 
database si trova, se il dato è aggiornato e di alta qualità, eliminando il rischio di errori o di ricreare set di dati che un collega aveva già preparato in passato.

\subsubsection{Data Quality e Governance}
Nel catalogo è possibile valutare la qualità delle informazioni, vedere chi è il responsabile di una specifica tabella e capire come i dati sono collegati tra loro. 
Assicura che chiunque in azienda prenda decisioni basandosi su informazioni affidabili, certificate e sicure, migliorando l'efficienza generale.


\section{Obiettivi dello stage}
% obiettivi dell'azienda (vedi prog. formativo) + altre cose tipo: studio completo del way of working aziendale per un mio possibile futuro in azienda
% studio delle tecnologie, dell'ambiente, confidenza con strumenti, competenze trasversali in più ambiti per tenere tutte le porte aperte

\subsection{Obiettivi personali}
% tipo aspettative personali e cosa mi aspetto/voglio dallo stage


\section{Vincoli di progetto}
\subsection{Vincoli tecnologici}
% per esempio le tecnologie che ho usato e che ho dovuto installare

\subsection{Vincoli metodologici}
% usare Teams per qualsiasi comubnicazione, studio della comunicazione (vedi doc), metodologia agile, aggironamenti frequenti, 
% scrittura di tutta la documentazione richiesta: dai documenti ufficiali per lo studio di OpenMetadata, ai verbali delle riunioni, a quelli relativi alla formazione, come
% i manuali di installazione o i verbali sugli incontri su Power BI

\subsection{Vincoli temporali}
% le ore a disposizione, considerando che il venerdì si esce un'ora prima e rispettando le pause. I vincoli temporali poi sono stati definiti nel piano di progetto che riporto dopo

\section{Pianificazione del lavoro}


\section{Formazione e tecnologie}
% formazione su SQL e database: query, sintassi SQL base, viste anche stored procedures, viste, cursori; confidenza con PostgreSQL e PgAdmin e anche con SQL Server; procedure di backup e restore
% docker e docker compose, WSL, VM (Virtual Box), perchè il pc è windows ma mi serviva anche ubuntu (+ studio comandi batch ubuntu)
% formazione su sviluppo sicuro, principi della privacy e protezione dei dati sensibili e personali
% altro: metodologia Agile (SCRUM), git e github, sistema gestionale di Wintech (ESOLVER), Power BI e principi della business intelligence




 %2
    % capitolo 3: analisi dei requisiti
\chapter{Analisi dei requisiti}
\label{cap:analisi-requisiti}
% TENERE SEMPRE A MENTE CHE OGNI REQUISITO DEVE RISPETTARE QUESTI 2 PRINCIPI:
% 1. Necessità: ogni requisito inserito nella lista deve soddisfare un bisogno aziendale reale ed esplicito.
% 2. Sufficienza: I requisiti elencati devono coprire tutti i bisogni dell'azienda.

 \section{Introduzione}
 In questa sezione vengono analizzati i casi d'uso e i requisiti aziendali necessari per l'introduzione di un \textit{Data Catalog} nell'infrastruttura di Wintech. 
 Dato che le moderne soluzioni software offrono un ampio spettro di funzionalità, non sempre necessarie, l'analisi ha lo scopo di filtrare gli strumenti attualmente 
 presenti sul mercato, individuando quelli che meglio rispondono ai requisiti dell'azienda e alle sue reali esigenze operative.

\subsection{Attori del sistema}
In ambito aziendale possono essere molteplici le figure che hanno la necessità di accedere a un \textit{Data Catalog}, ognuna con esigenze e competenze diverse. 
Ai fini dell'analisi dei requisiti, si individuano i seguenti attori principali:
\begin{itemize}
    \item \textbf{Analista o Sviluppatore di BI}: Utente finale del catalogo, il cui obiettivo è ricercare e accedere ai dati provenienti da varie fonti, comprenderne il significato e 
    utilizzarli per generare report e analisi o per esportarli in strumenti esterni di Business Intelligence.
    \item \textbf{Responsabile della Governance e della sicurezza dei dati}: Figura garante della qualità e della sicurezza dei dati, che controlla che i dati siano corretti e che rispettino le normative vigenti.
    \item \textbf{Amministratore di sistema}: Figura tecnica che gestisce l'infrastruttura della piattaforma e connette le varie fonti di dati aziendali e mantiene il catalogo aggiornato, garantendone il corretto funzionamento.
\end{itemize}

\newpage

\section{Casi d'uso}
In questa sezione vengono analizzati i casi d'uso e gli scenari per cui l'azienda ha bisogno di adottare un \textit{Data Catalog}. 
Per maggiore chiarezza organizzativa, i casi d'uso vengono suddivisi in tre macroaree raggruppate in base all'attore primario coinvolto.\\
Ogni caso d'uso è identificato dalla seguente notazione:
\begin{center}
    \textbf{UC-[Macroarea].[Numero]}
\end{center}

% MACROAREA 1
\subsection{Amministrazione del sistema}

\begin{usecase}{1.0}{Collegamento a una fonte dati strutturata}
    \usecaseactors{Amministratore di sistema}
    \usecasedesc{L'amministratore collega una fonte dati strutturata al catalogo dati aziendale e ne visualizza i metadati tramite interfaccia}
    \usecasepre{L'attore ha effettuato l'accesso alla piattaforma con privilegi di amministratore e possiede le credenziali della fonte dati da importare}
    \usecasepost{I metadati della fonte dati selezionata sono stati importati con successo e sono visibili all'interno del catalogo dati}
    \usecasemain{
        \begin{enumerate}
            \item L'attore seleziona l'opzione che permette di connettere una nuova fonte di dati al catalogo.
            \item L'attore seleziona la fonte desiderata e inserisce credenziali o parametri per la connessione.
            \item Il sistema verifica la connessione e notifica il successo della connessione.
            \item L'attore avvia l'estrazione dei metadati.
        \end{enumerate}
    }
    \usecasealt{
        \begin{enumerate}
            \item Il sistema segnala il fallimento della connessione.
            \item L'attore non può importare i dati nel catalogo e deve ricontrollare i parametri per la connessione.
        \end{enumerate}
    }
    \label{uc:1_0_fonte_strutturata}
\end{usecase}

\begin{usecase}{1.1}{Collegamento a una fonte dati semi strutturata}
    \usecaseactors{Amministratore di sistema}
    \usecasedesc{L'amministratore collega una fonte dati semi strutturata al catalogo dati aziendale e ne visualizza i metadati tramite interfaccia}
    \usecasepre{L'attore ha effettuato l'accesso alla piattaforma con privilegi di amministratore e possiede le credenziali o il percorso della fonte dati da importare}
    \usecasepost{I metadati della fonte dati selezionata sono stati importati con successo e sono visibili all'interno del catalogo dati}
    \usecasemain{
        \begin{enumerate}
            \item L'attore seleziona l'opzione che permette di aggiungere una nuova fonte di dati al catalogo.
            \item L'attore seleziona la fonte desiderata e inserisce i parametri per la connessione o il percorso del file.
            \item Il sistema verifica la connessione e notifica il successo della connessione.
            \item L'attore avvia l'estrazione dei metadati.
        \end{enumerate}
    }
    \label{uc:1_1_fonte_semi_strutturata}
\end{usecase}

\begin{usecase}{1.2}{Collegamento a una fonte dati non strutturata}
    \usecaseactors{Amministratore di sistema}
    \usecasedesc{L'amministratore collega una fonte dati non strutturata al catalogo dati aziendale e ne visualizza i metadati tramite l'interfaccia della piattaforma}
    \usecasepre{L'attore ha effettuato l'accesso alla piattaforma con privilegi di amministratore e possiede le credenziali o il percorso della fonte dati da importare}
    \usecasepost{I metadati della fonte dati selezionata sono stati importati con successo e sono visibili all'interno del catalogo dati}
    \usecasemain{
        \begin{enumerate}
            \item L'attore seleziona l'opzione che permette di aggiungere una nuova fonte di dati al catalogo.
            \item L'attore seleziona la fonte desiderata e inserisce i parametri per la connessione o il percorso del file.
            \item Il sistema verifica la connessione e notifica il successo della connessione.
            \item L'attore avvia l'estrazione dei metadati.
        \end{enumerate}
    }
    \label{uc:1_2_fonte_non_strutturata}
\end{usecase}

\begin{usecase}{1.3}{Visualizzazione delle fonti dati collegate}
    \usecaseactors{Amministratore di sistema}
    \usecasedesc{L'amministratore visualizza l'inventario delle fonti dati, cioè l'insieme delle fonti di dati che popolano il \textit{Data Catalog}}
    \usecasepre{L'attore ha effettuato l'accesso alla piattaforma con privilegi di amministratore}
    \usecasepost{L'attore accede all'elenco completo delle fonti di dati importate nel \textit{Data Catalog}}
    \usecasemain{
        \begin{enumerate}
            \item L'attore seleziona l'opzione per visualizzare l'elenco delle fonti dati del catalogo dati.
            \item Il sistema mostra l'elenco completo delle fonti dati.
        \end{enumerate}
    }
    \label{uc:1_3_visualizzazione_fonti_collegate}
\end{usecase}

\begin{usecase}{1.4}{Gestione dei permessi e degli accessi}
    \usecaseactors{Amministratore di sistema}
    \usecasedesc{L'amministratore definisce i ruoli degli utenti che accedono alla piattaforma e assegna i permessi di accesso ai diversi set di dati. 
    Il sistema garantisce che ogni utente visualizzi esclusivamente i dati per cui dispone delle autorizzazioni necessarie}    
    \usecasepre{L'attore ha effettuato l'accesso alla piattaforma con privilegi di amministratore}
    \usecasepost{I ruoli degli utenti sono stati definiti e i permessi di accesso ai dati sono stati configurati correttamente}
    \usecasemain{
        \begin{enumerate}
            \item L'attore seleziona l'opzione per la gestione dei ruoli e dei permessi.
            \item L'attore crea o modifica i ruoli disponibili nella piattaforma.
            \item L'attore assegna i ruoli agli utenti che accedono al catalogo.
            \item L'attore definisce le regole di accesso ai dati associate a ciascun ruolo.
            \item Il sistema applica i permessi configurati e limita la visibilità dei dati in base al ruolo di ciascun utente.
        \end{enumerate}
    }
    \label{uc:1_4_gestione_permessi}
\end{usecase}

\begin{usecase}{1.5}{Pianificazione della sincronizzazione dei metadati}
    \usecaseactors{Amministratore di sistema}
    \usecasedesc{L'amministratore pianifica la sincronizzazione automatica dei metadati, definendo la frequenza e le modalità con cui il catalogo si aggiorna allineandosi 
    alle fonti dati originali, garantendo che le informazioni presenti nel catalogo siano sempre coerenti e aggiornate}
    \usecasepre{L'attore ha effettuato l'accesso alla piattaforma con privilegi di amministratore e almeno una fonte dati è stata collegata al catalogo}
    \usecasepost{La sincronizzazione automatica è stata pianificata correttamente e i metadati vengono aggiornati secondo la frequenza stabilita}
    \usecasemain{
        \begin{enumerate}
            \item L'attore seleziona le impostazioni per la pianificazione della sincronizzazione dei metadati.
            \item L'attore configura la frequenza e le modalità di aggiornamento desiderate.
            \item Il sistema conferma la configurazione e pianifica le sincronizzazioni secondo i parametri definiti.
            \item Il sistema esegue automaticamente la sincronizzazione dei metadati agli intervalli stabiliti.
        \end{enumerate}
    }
    \label{uc:1_5_pianificazione_sincronizzazione}
\end{usecase}

\begin{usecase}{1.6}{Configurazione della classificazione automatica dei dati sensibili}
    \usecaseactors{Amministratore di sistema}
    \usecasedesc{L'amministratore abilita e configura il motore di classificazione automatica dei dati sensibili, che analizza i metadati delle fonti dati collegate 
    al catalogo e individua in autonomia i campi contenenti dati personali o sensibili (PII). Il motore si avvale per esempio di tecniche di Intelligenza Artificiale e di espressioni 
    regolari per riconoscere pattern tipici di dati sensibili, quali nomi, indirizzi email, numeri di telefono, codici fiscali e altri dati identificativi}
    \usecasepre{L'attore ha effettuato l'accesso alla piattaforma con privilegi di amministratore e almeno una fonte dati è stata collegata al catalogo}
    \usecasepost{Il motore di classificazione automatica è attivo e configurato. Durante le sincronizzazioni dei metadati, il sistema individua e contrassegna automaticamente i campi contenenti dati sensibili}
    \usecasemain{
        \begin{enumerate}
            \item L'attore accede alle impostazioni della classificazione automatica dei dati sensibili.
            \item L'attore abilita il motore di classificazione e seleziona le fonti dati su cui applicarlo.
            \item L'attore configura le categorie di dati sensibili da rilevare o adotta la configurazione predefinita della piattaforma.
            \item Il sistema conferma la configurazione e attiva il motore.
            \item Ad ogni sincronizzazione, il sistema analizza i metadati delle fonti selezionate e contrassegna automaticamente i campi identificati come sensibili.
        \end{enumerate}
    }
    \label{uc:1_6_configurazione_classificazione_dati}
\end{usecase}

\begin{usecase}{1.7}{Monitoraggio dello stato delle sincronizzazioni e delle connessioni}
    \usecaseactors{Amministratore di sistema}
    \usecasedesc{L'amministratore visualizza una dashboard tecnica, un report o un avviso che indica lo stato delle connessioni alle fonti dati e l'esito 
    delle sincronizzazioni dei metadati, consentendo di verificare che le operazioni pianificate siano avvenute correttamente e di individuare eventuali anomalie o errori}
    \usecasepre{L'attore ha effettuato l'accesso alla piattaforma con privilegi di amministratore e almeno una sincronizzazione è stata pianificata ed eseguita}
    \usecasepost{L'attore ha visualizzato lo stato aggiornato delle connessioni e delle sincronizzazioni ed è a conoscenza di eventuali errori o anomalie da risolvere}
    \usecasemain{
        \begin{enumerate}
            \item L'amministratore accede alla dashboard di monitoraggio delle connessioni e delle sincronizzazioni.
            \item Il sistema mostra lo stato di ciascuna operazione, indicando se è avvenuta con successo o se si sono verificati errori.
            \item In caso di errore, il sistema fornisce informazioni dettagliate sulla causa del problema, ad esempio una credenziale scaduta o una fonte dati non raggiungibile.
            \item L'attore prende visione delle anomalie e interviene per risolvere il problema alla fonte.
        \end{enumerate}
    }
    \label{uc:1_7_monitoraggio_sincronizzazioni}
\end{usecase}

\begin{usecase}{1.8}{Gestione di avvisi e alert tecnici}
    \usecaseactors{Amministratore di sistema}
    \usecasedesc{Il sistema invia automaticamente alert di natura tecnica e avvisi all'amministratore quando vengono rilevate anomalie tecniche, come il fallimento di una sincronizzazione, 
    un errore di connessione o la scadenza delle credenziali. L'amministratore visualizza gli alert ricevuti e interviene per risolvere i problemi segnalati}
    \usecasepre{L'attore ha effettuato l'accesso alla piattaforma con privilegi di amministratore e il sistema di notifiche automatiche è stato configurato}
    \usecasepost{L'attore ha preso visione degli alert ricevuti e ha avviato le azioni correttive necessarie per risolvere i problemi segnalati}
    \usecasemain{
        \begin{enumerate}
            \item Il sistema rileva un'anomalia e genera automaticamente un alert.
            \item L'attore riceve una notifica e accede alla sezione dedicata agli alert e agli avvisi.
            \item Il sistema mostra i dettagli dell'alert, come la causa, la fonte dati coinvolta e il momento in cui si è verificato il problema.
            \item L'attore esamina le informazioni disponibili e interviene per risolvere il problema.
            \item Il sistema non mostra più l'alert o lo marchia come "risolto".
        \end{enumerate}
    }
    \label{uc:1_8_gestione_alert_tecnici}
\end{usecase}


% MACROAREA 2
\subsection{Governance e sicurezza dei dati}
\begin{usecase}{2.0}{Individuazione automatica dei dati sensibili}
    \usecaseactors{Responsabile della Governance e della sicurezza dei dati}
    \usecasedesc{La piattaforma analizza automaticamente i metadati dei dataset caricati nel catalogo e individua i campi contenenti dati sensibili o personali (PII), classificandoli tramite tag appositi. 
    Il processo avviene in modo trasparente durante la sincronizzazione dei metadati, senza richiedere intervento manuale da parte dell'attore}
    \usecasepre{L'attore ha effettuato l'accesso alla piattaforma con privilegi di responsabile della Governance e il motore di classificazione automatica è già stato configurato dall'amministratore}
    \usecasepost{I campi contenenti dati sensibili o personali sono stati identificati e classificati automaticamente con i tag appropriati}
    \usecasemain{
        \begin{enumerate}
            \item La piattaforma rileva l'aggiunta o la sincronizzazione di un nuovo dataset nel catalogo.
            \item Il motore di classificazione automatica analizza i metadati del dataset e individua i campi che corrispondono a pattern tipici di dati sensibili o personali.
            \item Il sistema associa automaticamente i tag appropriati ai campi individuati.
            \item L'attore riceve una notifica dell'avvenuta classificazione automatica e può procedere alla revisione o all'approvazione.
        \end{enumerate}
    }
    \label{uc:2_0_individuazione_automatica_dati}
\end{usecase}




% MACROAREA 3
\subsection{Analisi dei dati e Business Intelligence}


\newpage

\section{Requisiti}
In questa sezione vengono elencati e descritti i requisiti derivanti dallo studio dei casi d'uso, con l'obiettivo di avere un quadro chiaro delle caratteristiche 
che un \textit{Data Catalog} deve avere per rispondere alle esigenze aziendali. Per ogni requisito vengono forniti una descrizione e il tracciamento verso i casi d'uso correlati.\\

Ogni requisito è identificato dalla seguente notazione:
\begin{center}
    \textbf{R[Tipo]-[Importanza]-[Numero]}
\end{center}
dove le voci assumono i seguenti significati:
\begin{description}
    \item[Tipo] Indica la natura architetturale del requisito:
    \begin{itemize}
        \item \textbf{F (Funzionale)}: Descrive i servizi e le funzionalità tangibili che il sistema deve fornire.
        \item \textbf{V (Vincolo)}: Descrive limitazioni tecnologiche o normative imposte esternamente.
        \item \textbf{Q (Qualità)}: Descrive le caratteristiche qualitative misurabili del sistema.
    \end{itemize}

    \item[Importanza] Indica il grado di priorità per l'azienda:
    \begin{itemize}
        \item \textbf{O (Obbligatorio)}: Requisito irrinunciabile; la sua assenza compromette gli obiettivi primari del progetto.
        \item \textbf{D (Desiderabile)}: Requisito non strettamente vincolante nel breve termine, ma a valore aggiunto riconoscibile.
        \item \textbf{OP (Opzionale)}: Requisito di priorità minore, la cui implementazione può essere rimandata.
    \end{itemize}

    \item[Numero] Un intero progressivo univoco per ogni tipologia.
\end{description}

% TABELLE DEI REQUISITI
\subsection{Requisiti funzionali}
I requisiti funzionali descrivono i comportamenti e i servizi che il sistema deve fornire. Rispondono alla domanda: \textit{"Cosa deve fare il software?"}.


\subsection{Requisiti di vincolo}
I requisiti di vincolo rappresentano le limitazioni tecnologiche, normative, architetturali o implementative imposte dall'esterno entro le quali il progetto deve muoversi.

\subsection{Requisiti di qualità}
I requisiti di qualità descrivono le caratteristiche misurabili che il sistema deve garantire, come prestazioni, sicurezza, usabilità o manutenibilità.




\newpage

\section{Tracciamento dei requisiti}
% tabella di tracciamento fonte - requisito




% ROBA DEL TEMPLATE ==================================

% \begin{usecase}{0}{Scenario principale}
% \usecaseactors{Sviluppatore applicativi}
% \usecasepre{Lo sviluppatore è entrato nel plug-in di simulazione all'interno dell'IDE}
% \usecasedesc{La finestra di simulazione mette a disposizione i comandi per configurare, registrare o eseguire un test}
% \usecasepost{Il sistema è pronto per permettere una nuova interazione}
% \label{uc:scenario-principale}
% \end{usecase}


% \begin{table}%
% \caption{Tabella del tracciamento dei requisti funzionali}
% \label{tab:requisiti-funzionali}
% \begin{tabularx}{\textwidth}{lXl}
% \hline\hline
% \textbf{Requisito} & \textbf{Descrizione} & \textbf{Use Case}\\
% \hline
% RFN-1     & L'interfaccia permette di configurare il tipo di sonde del test & UC1 \\
% \hline
% \end{tabularx}
% \end{table}
 %3
    % capitolo 4: valutazione di OpenMetadata e dei competitor

\chapter{Valutazione di OpenMetadata e dei Competitor}
\label{cap:valutazione-competitor}

\intro{Breve introduzione al capitolo}\\

 %4
    % capitolo 5 - Installazione e configurazione + data ingestion dei primi dataset
\chapter{Installazione e configurazione}
\label{cap:installazione}

\intro{In questo capitolo vengono forniti i prerequisiti e le indicazioni per l'installazione e il \emph{setup} di OpenMetadata in ambiente locale.}\\

\section{Prerequisiti}
L'installazione e l'esecuzione di OpenMetadata prevedono requisiti specifici, legati alle risorse hardware della macchina ospitante e 
all'ambiente di esecuzione locale, basato in questo contesto sulla containerizzazione tramite Docker.

Per garantire il corretto funzionamento di una distribuzione locale di test, la macchina ospitante (basata su sistema operativo Linux oppure su Windows 
tramite \gls{wsl}2 o macchina virtuale) deve allocare all'\emph{engine} di Docker le seguenti risorse hardware minime:
\begin{itemize}
    \item \textbf{CPU}: Almeno 4 vCPU.
    \item \textbf{Memoria RAM}: Almeno 6--8 GB di RAM interamente dedicati ai container.
    \item \textbf{Spazio su disco}: Almeno 30--40 GB di spazio libero.
\end{itemize}

Dal punto di vista software, i requisiti fondamentali prevedono Git per il controllo di versione, Docker Engine (versione 20.10.0 o superiore) e Docker Compose (versione v2.1.1 o superiore). 
L'architettura locale di OpenMetadata viene orchestrata attraverso quattro container distinti:
\begin{enumerate}
    \item \textbf{Database di sistema}: Un'istanza PostgreSQL o MySQL dedicata al salvataggio dello stato, delle configurazioni e dei metadati gestiti dalla piattaforma.
    \item \textbf{OpenMetadata Server}: Il backend centrale della piattaforma sviluppato in Java.
    \item \textbf{OpenMetadata Ingestion (\emph{Apache Airflow})}: Il motore di orchestrazione responsabile dell'esecuzione automatica e pianificata delle pipeline di estrazione dati.
    \item \textbf{Motore di ricerca}: Un container ElasticSearch (o \gls{opensearch}) dedicato all'indicizzazione dei metadati e alla ricerca \emph{full-text}.
\end{enumerate}

Per consentire la corretta fruizione dei servizi sul browser, l'ambiente di rete locale deve garantire che le porte 8585 (dedicata all'interfaccia utente di OpenMetadata) 
e 8080 (dedicata al pannello di controllo di Apache Airflow) siano libere da conflitti.

Infine, per testare la piattaforma senza violare i vincoli di riservatezza sui dati reali aziendali, durante lo stage è stato utilizzato il 
database relazionale di esempio \emph{AdventureWorks} fornito da Microsoft\cite{site:adventureworks}. 
Tale risorsa è stata ripristinata all'interno di un'istanza locale di Microsoft SQL Server, simulando uno scenario di produzione realistico, completo di 
relazioni e chiavi esterne su cui eseguire le prime procedure di \emph{Data Ingestion} e \emph{Data Discovery}.

\section{Installazione di OpenMetadata}
Una volta soddisfatti tutti i prerequisiti, la procedura di deployment locale di OpenMetadata è stata eseguita tramite Docker e Docker Compose, 
seguendo i passaggi della documentazione ufficiale\cite{site:openmetadata-docs}, di seguito elencati:
\begin{enumerate}
    \item \textbf{Creazione della directory di lavoro}: è stata creata una cartella dedicata al progetto:
    \begin{verbatim}
mkdir openmetadata-docker && cd openmetadata-docker
    \end{verbatim}
    \item \textbf{Download del file di configurazione}: è stato scaricato il file \texttt{docker-compose.yml} ufficiale dal repository\cite{site:openmetadata-repo} del progetto (oppure il file \texttt{docker-compose-postgres.yml} qualora si preferisca utilizzare PostgreSQL come database di sistema al posto di MySQL).
    \item \textbf{Avvio dei servizi}: i container Docker di OpenMetadata sono stati eseguiti tramite il seguente comando:
    \begin{verbatim}
docker compose -f docker-compose.yml up --detach
    \end{verbatim}
\end{enumerate}

Una volta che i container risultano operativi, la piattaforma è pronta per l'utilizzo:
\begin{itemize}
    \item \textbf{Interfaccia utente di OpenMetadata}: raggiungibile da browser all'indirizzo \nolinkurl{http://localhost:8585}. Al primo accesso, l'autenticazione avviene tramite le credenziali di default fornite dal sistema: \texttt{admin@open-metadata.org} come \emph{username} e \texttt{admin} come \emph{password}. Da qui è poi possibile accedere alle impostazioni generali e configurare nuovi utenti;
    \item \textbf{Pannello di controllo di Apache Airflow}: raggiungibile all'indirizzo \nolinkurl{http://localhost:8080} inserendo le credenziali di default (\texttt{admin} / \texttt{admin}). Airflow opera in background come motore di orchestrazione per le pipeline di \emph{ingestion} dei metadati; sebbene la gestione delle sincronizzazioni avvenga interamente dall'interfaccia di OpenMetadata, l'accesso diretto ad Airflow risulta fondamentale per l'analisi dei log e il \emph{debugging} avanzato in caso di fallimento delle pipeline.
\end{itemize}


\section{Data Ingestion}

\section{Ricerca e Data Discovery}

\section{Data Lineage} %5
    % capitolo 6: data quality e observability
\chapter{Data Quality e Data Observability}
\label{cap:quality-observability}

\intro{Il seguente capitolo illustra le funzionalità di Data Quality e gli strumenti di monitoraggio dei dati nativi di OpenMetadata.}\\

\section{Agenti automatici}
All'interno di OpenMetadata, la raccolta continuativa dei metadati e il monitoraggio dello stato di salute delle informazioni sono affidati al motore di automazione \textbf{AutoPilot}. 
Si tratta di un sistema di agenti automatici che orchestrano le operazioni di estrazione, profilazione, classificazione e tracciamento senza richiedere configurazioni 
manuali complesse o interventi ripetitivi da parte dell'utente.

Gli agenti si attivano automaticamente nel momento in cui viene registrato un nuovo servizio o quando viene avviata una sincronizzazione con le fonti originali, 
operando in background secondo una pianificazione personalizzabile. Ciascun agente assolve a compiti specifici e mirati all'interno della piattaforma:
\begin{itemize}
    \item \textbf{Metadata Agent}: Costituisce l'agente primario ed è prerequisito per il funzionamento dell'intera piattaforma. Si occupa dell'estrazione e della sincronizzazione dei metadati dalle sorgenti collegate.
    \item \textbf{Profiler Agent}: Raccoglie metriche e statistiche descrittive sui dati a livello di tabella e di singola colonna (per esempio, calcolando la percentuale di valori nulli, il numero di valori distinti, gli intervalli minimi e massimi). I risultati alimentano direttamente la scheda di \emph{Data Quality} di ciascun dataset e forniscono la base quantitativa su cui vengono valutati i test di qualità.
    \item \textbf{AutoClassification Agent}: Analizza i campioni di testo prelevati dal Profiler Agent per alimentare il motore di \emph{Entity Recognition} basato su tecniche NLP, identificando automaticamente le informazioni personali e sensibili e applicando i relativi tag di sicurezza.
    \item \textbf{Lineage Agent}: Analizza i log delle query SQL e le definizioni delle viste per ricostruire in modo automatico il grafo delle dipendenze tra i dati.
    \item \textbf{Usage Agent e agenti specializzati}: Raccolgono le statistiche di utilizzo delle tabelle per alimentare il punteggio di popolarità (\emph{Popularity Scoring}) e le \emph{Query Analytics}.
\end{itemize}

Ogni agente è responsabile di una fase specifica del ciclo di vita dei metadati e opera in sequenza logica: poiché l'esecuzione di alcuni agenti è prerequisito per l'attivazione di altri, 
è fondamentale che vengano eseguiti nell'ordine corretto.

Lo stato di ciascun agente è visibile all'interno della scheda \emph{Agents} nella pagina del servizio di riferimento. Da questa interfaccia, gli amministratori e gli utenti 
autorizzati possono controllare lo stato di avanzamento delle attività (\emph{Active}, \emph{Pending}, \emph{Success}, \emph{Failed}), ispezionare i log di esecuzione per il 
\emph{debugging} o forzare un'esecuzione manuale a seguito di modifiche strutturali apportate a uno specifico dataset o servizio sorgente.


\section{Test di qualità}

\section{Gestione degli incidenti}

\section{Sistemi di alerting e notifica}
% e anche data collaboration %6
    % capitolo 7: data governance e sviluppo sicuro
\chapter{Data Governance e Sviluppo Sicuro}
\label{cap:sviluppo-sicuro}

\intro{Questo capitolo affronta le tematiche legate alla sicurezza e alla conformità normativa, in linea con i principi della \emph{Privacy by Design}.}\\

\section{Introduzione allo sviluppo sicuro}
L'introduzione di un Data Catalog all'interno dell'infrastruttura aziendale non risponde solo alle esigenze di catalogazione dei dati e di gestione dei metadati, 
ma rappresenta anche uno strumento per garantire la conformità normativa dei dati, l'adesione alle policy aziendali e il rispetto del \gls{gdpr}, nonché un mezzo per supportare il Ciclo di Sviluppo del Software Sicuro.

A tal proposito, Wintech mette a disposizione una serie di documenti e linee guida per la protezione della privacy, la gestione dei dati sensibili e lo sviluppo sicuro del software.

All'interno del contesto aziendale, viene definito il concetto di \emph{Privacy by Design} (PbD), intesa come un approccio olistico concettuale che può essere applicato \emph{end-to-end} all'interno di un'organizzazione, 
includendo le sue tecnologie informatiche, le sue pratiche commerciali, i suoi processi, la progettazione fisica e le infrastrutture di rete. 
Questa definizione si basa su tre concetti fondamentali:
\begin{itemize}
    \item \textbf{Integrazione nativa fin dalla progettazione}: Qualsiasi progetto deve essere realizzato considerando la riservatezza e la protezione dei dati personali sin dalla sua prima fase di progettazione e nell'architettura dei sistemi IT. Non è una misura attuata successivamente a un evento o a un problema sorto.
    \item \textbf{Approccio \emph{user centric}}: Secondo questa impostazione, l'utente deve essere posto al centro del sistema di protezione dei dati personali.
    \item \textbf{Azione proattiva e preventiva}: La PbD è caratterizzata dall'adozione di misure proattive piuttosto che reattive. Il suo scopo primario è anticipare e prevenire gli eventi lesivi per la privacy prima che accadano, invece di attendere che il rischio si materializzi per poi correggerlo.
\end{itemize}

\subsection{Principi della privacy}
Nella definizione delle proprie Linee Guida allo Sviluppo Sicuro, Wintech si ispira allo standard \textbf{ISO/IEC 29100:2011}, il quale raccoglie undici principi fondamentali 
che indirizzano la progettazione, lo sviluppo e l'implementazione dei requisiti di protezione della privacy all'interno dei software. 
Tali principi sono descritti in dettaglio di seguito:

% \begin{longtable}{ p{0.30\textwidth} p{0.64\textwidth} }
\begin{longtable}{ >{\raggedright\arraybackslash}p{0.30\textwidth} p{0.64\textwidth} }
\toprule
\textbf{Principio} & \textbf{Descrizione} \\
\midrule
\endfirsthead

\toprule
\textbf{Principio} & \textbf{Descrizione} \\
\midrule
\endhead

\midrule
\noalign{\vspace{10pt}}
\caption[]{Principi della privacy secondo lo standard ISO/IEC 29100:2011} \\
\endfoot

\bottomrule
\noalign{\vspace{10pt}}
\caption{Principi della privacy secondo lo standard ISO/IEC 29100:2011}
\label{tab:principi-privacy}
\endlastfoot
    \textbf{Consenso e Scelta} & L'interessato deve poter scegliere se acconsentire o meno al trattamento, in modo chiaro e accessibile.\\ \midrule
    \textbf{Legittimità e Specifica delle Finalità} & Il trattamento deve essere conforme alla legge e basato su una base giuridica ammissibile.\\ \midrule
    \textbf{Limitazione della Raccolta} & La raccolta dei dati deve limitarsi a ciò che è strettamente necessario per gli scopi prefissati.\\ \midrule
    \textbf{Minimizzazione dei Dati} & I sistemi ICT devono essere progettati per ridurre al minimo i dati personali elaborati.\\ \midrule
    \textbf{Limitazione dell'Utilizzo, Conservazione e Divulgazione} & I dati vanno usati e conservati solo per scopi specifici, espliciti e legittimi.\\ \midrule
    \textbf{Accuratezza e Qualità} & I dati elaborati devono essere accurati, completi, aggiornati e pertinenti.\\ \midrule
    \textbf{Apertura, Trasparenza e Informativa} & Devono essere fornite informazioni chiare e accessibili sulle procedure di trattamento.\\ \midrule
    \textbf{Partecipazione Individuale e Accesso} & Agli interessati deve essere garantita la possibilità di accedere e rivedere i propri dati.\\ \midrule
    \textbf{Responsabilizzazione} & Le politiche di privacy devono essere documentate, comunicate e assegnate a un responsabile specifico.\\ \midrule
    \textbf{Sicurezza delle Informazioni} & I dati devono essere protetti da rischi come accesso non autorizzato, modifica, divulgazione o perdita.\\ \midrule
    \textbf{Conformità alla Privacy} & Devono essere previste verifiche periodiche per dimostrare il rispetto della tutela della privacy.\\
\end{longtable}

\subsection{Adeguamento al GDPR nel Ciclo di Sviluppo del Software Sicuro}
Il \gls{gdpr} ha un impatto profondo sul \gls{sdlc} per tutte le organizzazioni che progettano e implementano sistemi informativi all'interno dell'Unione Europea. 
I requisiti funzionali e tecnici introdotti dal Regolamento influenzano quasi tutti gli aspetti della progettazione architetturale, a partire 
dai livelli più bassi, come l'architettura dei dati (database) e i livelli di trasporto, per arrivare fino ai livelli applicativi e di presentazione. 
Wintech si impegna a soddisfare tali requisiti fin dalle prime fasi dell'SDLC.

Di seguito vengono mappati e spiegati alcuni dei principali articoli del GDPR e le misure che l'azienda si impegna ad adottare per garantire la \emph{Privacy by Design} e lo Sviluppo Sicuro.
% \begin{longtable}{ p{0.18\textwidth} p{0.37\textwidth} p{0.37\textwidth} }
\begin{longtable}{ >{\raggedright\arraybackslash}p{0.18\textwidth} p{0.37\textwidth} p{0.37\textwidth} }
\toprule
\textbf{Articolo del GDPR} & \textbf{Principio normativo} & \textbf{Impegno Wintech} \\
\midrule
\endfirsthead

\toprule
\textbf{Articolo del GDPR} & \textbf{Principio normativo} & \textbf{Impegno Wintech} \\
\midrule
\endhead

\midrule
\noalign{\vspace{10pt}}
\caption[]{Mappatura dei principali articoli del GDPR rispetto agli impegni aziendali} \\
\endfoot

\bottomrule
\noalign{\vspace{10pt}}
\caption{Mappatura dei principali articoli del GDPR rispetto agli impegni aziendali}
\label{tab:articoli-gdpr}
\endlastfoot
    \textbf{Art. 5 -- Principi applicabili al trattamento dei dati personali} & I dati personali sono raccolti in modo lecito, corretto e trasparente nei confronti dell'interessato e per finalità determinate, esplicite e legittime. Devono essere adeguati, pertinenti e limitati a quanto necessario per le finalità dichiarate; devono essere esatti, aggiornati e conservati per un arco di tempo non superiore al conseguimento delle finalità stesse. & Sotto la direzione del titolare, i dati sono protetti garantendone l'integrità e la riservatezza. Wintech implementa strategie di gestione dei registri elettronici affinché i dati vengano conservati solo per il tempo strettamente necessario; i dati personali scaduti non vengono memorizzati e l'interessato viene informato di tale periodo.\\ \midrule
    \textbf{Art. 13 -- Informazioni da fornire al momento della raccolta} & Il titolare del trattamento deve fornire all'interessato le proprie informazioni di contatto, le finalità e la base giuridica del trattamento, gli eventuali destinatari, il periodo di conservazione, nonché l'esistenza dei diritti di accesso, rettifica, cancellazione e revoca del consenso. & I sistemi sono progettati affinché, al momento della raccolta dei dati, gli attributi del trattamento e le fasi elaborative vengano forniti all'utente per via elettronica o per iscritto, in forma chiara e facilmente comprensibile, garantendo assoluta trasparenza.\\ \midrule
    \textbf{Art. 15 -- Diritto di accesso dell'interessato} & L'interessato ha il diritto di ottenere dal titolare del trattamento la conferma che sia o meno in corso un trattamento di dati personali che lo riguardano e, in tal caso, di ottenerne l'accesso. & Le architetture software devono garantire all'interessato il diritto di accedere ai propri dati e di controllarne il trattamento in qualsiasi momento. Inoltre, l'infrastruttura deve prevedere la possibilità di fornire all'interessato una copia dei propri dati in un formato comunemente utilizzato.\\ \midrule
    \textbf{Art. 16 -- Diritto di rettifica} & L'interessato ha il diritto di ottenere dal titolare del trattamento la rettifica dei dati personali inesatti che lo riguardano senza ingiustificato ritardo e l'integrazione dei dati incompleti. & I sistemi informativi devono consentire all'interessato di ottenere gratuitamente l'aggiornamento dei propri dati in caso di errore. L'infrastruttura deve supportare l'esercizio di tale diritto richiedendo l'aggiornamento direttamente per via elettronica.\\ \midrule
    \textbf{Art. 17 -- Diritto alla cancellazione (Diritto all'oblio)} & L'interessato ha il diritto di ottenere dal titolare la cancellazione dei dati personali che lo riguardano senza ingiustificato ritardo, qualora sussistano i motivi previsti dal Regolamento (es. revoca del consenso o cessazione delle finalità). & I sistemi informativi devono permettere la cancellazione immediata dei dati qualora richiesta. Inoltre, il dipartimento IT si impegna a mantenere inventari aggiornati dei dati interni ed esterni, in modo da sapere esattamente dove tali dati sono conservati da terze parti e poter comunicare tempestivamente l'ordine di cancellazione ad altre organizzazioni IT.\\ \midrule
    \textbf{Art. 25 -- Protezione dei dati fin dalla progettazione} & Il titolare del trattamento mette in atto misure tecniche e organizzative adeguate a garantire che siano trattati solo i dati personali necessari per ogni specifica finalità, integrandovi le necessarie garanzie a tutela dei diritti degli interessati. & L'implementazione della protezione dei dati all'interno del sistema informativo e dell'organizzazione avviene in modo nativo, per progettazione e come impostazione predefinita. Wintech considera questo approccio un requisito necessario nello sviluppo IT.\\ \midrule
    \textbf{Art. 32 -- Sicurezza del trattamento} & Il titolare e il responsabile del trattamento mettono in atto misure tecniche e organizzative per garantire un livello di sicurezza adeguato al rischio, tra cui la pseudonimizzazione e la cifratura dei dati personali. & L'infrastruttura IT di Wintech deve essere progettata avvalendosi di mezzi tecnici e organizzativi rigorosi. Ove possibile, i sistemi software sviluppati devono prevedere sempre l'uso della cifratura dei dati, nonché supportare l'anonimizzazione o la pseudonimizzazione delle informazioni.\\ \midrule
    \textbf{Artt. 33 e 34 -- Notifica e comunicazione di violazione dei dati} & In caso di violazione della sicurezza che comporta la perdita o la divulgazione illecita di dati, il titolare deve notificarla all'autorità di controllo (entro 72 ore) e all'interessato, qualora la violazione presenti un rischio elevato per i suoi diritti. & I sistemi software devono essere dotati di un ambiente tecnico opportuno per individuare, tracciare e valutare quasi immediatamente le violazioni dei dati. Questo è necessario per stabilire in tempo reale se vi sia un rischio elevato e adempiere agli obblighi di comunicazione.\\
\end{longtable}

La gestione di tutti i requisiti normativi sopraelencati richiede una documentazione costante delle fasi di progettazione e un supporto tecnologico adeguato. 
In questo scenario, l'adozione di un Data Catalog aziendale come OpenMetadata diventa una componente architetturale molto importante: automatizzando l'individuazione dei dati sensibili, 
applicando i relativi tag di sicurezza e mappando i flussi informativi, il catalogo fornisce all'azienda uno strumento concreto per garantire la conformità al ciclo di sviluppo sicuro  
e alle direttive imposte dal GDPR, trasformando i vincoli legali in regole tecniche implementate in modo proattivo.


\section{Classificazione dei dati}
Nel contesto della Data Governance, la classificazione dei dati rappresenta il processo logico mediante il quale si associano etichette semantiche (tag) agli asset informativi 
per descriverne in modo standardizzato la natura, l'area di business, l'appartenenza aziendale e il livello di riservatezza. Questa operazione agevola le attività di Data Discovery, 
consentendo agli utenti di cercare e filtrare gli elementi del catalogo, e sta alla base del rispetto delle policy di sicurezza e della conformità normativa, con particolare riferimento al GDPR.

Per configurare correttamente un catalogo dati aziendale, è fondamentale riconoscere quali tipologie di informazioni necessitano di maggiore tutela, come i 
dati sensibili e personali. A tal proposito, l'Articolo 4 del GDPR fornisce la seguente definizione di \textbf{dato personale}~\cite{site:gdpr}:
\begin{quoting}
    «Qualsiasi informazione riguardante una persona fisica identificata o identificabile (\emph{interessato}); si considera identificabile la persona fisica che può essere identificata, direttamente o indirettamente, con particolare riferimento a un identificativo come 
    il nome, un numero di identificazione, dati relativi all'ubicazione, un identificativo online o a uno o più elementi caratteristici della sua identità fisica, fisiologica, genetica, psichica, economica, culturale o sociale.»
\end{quoting}

Le linee guida aziendali di Wintech si adeguano a questa normativa, classificando le informazioni nelle seguenti macro-categorie:
\begin{itemize}
    \item \textbf{Dati identificativi (\gls{pii})}: Dati che consentono di identificare direttamente o indirettamente una persona.
    \item \textbf{Dati particolari/sensibili (Art. 9 GDPR)}: Dati personali che necessitano di una maggiore protezione poiché il loro trattamento potrebbe creare rischi significativi ai diritti o alle libertà fondamentali dell'interessato; si tratta per esempio di: origine razziale o etnica, opinione politica, ideologie religiose o filosofiche, appartenenza sindacale, dati genetici o biometrici, dati relativi alla salute o alla vita sessuale.
    \item \textbf{Dati giudiziari}: Dati relativi alle condanne penali e ai reati o a connesse misure di sicurezza; la gestione di queste informazioni è consentita sotto il controllo di un'autorità pubblica.
\end{itemize}

\noindent Viene fatta anche una distinzione tra i seguenti concetti:
\begin{itemize}
    \item \textbf{Dati pseudonimizzati}: Dati che vengono sostituiti da pseudonimi, in cui l'assegnazione degli pseudonimi è tale da non poter essere invertita e gli attributi collegati allo pseudonimo non sono sufficienti per identificare l'interessato. Tuttavia, i dati pseudonimizzati possono essere violati e l'interessato può essere identificato, pertanto sono soggetti al GDPR.
    \item \textbf{Dati anonimizzati}: Informazioni modificate in modo irreversibile, al punto che il titolare del trattamento non può identificare l'interessato. I dati resi anonimi non sono più considerati personali e quindi escono dal perimetro del GDPR.
    \item \textbf{Metadati}: I dati personali possono essere memorizzati in un sistema ICT in modo tale da non essere facilmente visibili all'utente del sistema. L'interessato deve essere a conoscenza dell'esistenza di tali dati sotto forma di metadati o del loro trattamento per tale scopo.
    \item \textbf{Dati non richiesti}: Informazioni personali fornite volontariamente dall'utente anche quando non richieste dal trattamento. Il rischio di raccogliere informazioni personali indesiderate può essere ridotto considerando le misure di tutela della privacy al momento della progettazione del software.
\end{itemize}

\noindent Data la specificità di queste categorie, una gestione manuale efficace della classificazione e del controllo dei dati personali risulterebbe molto complessa. 
In questo contesto, la funzionalità di auto-classificazione di un Data Catalog come OpenMetadata assume un ruolo chiave: la piattaforma scansiona gli asset di dati 
alla ricerca di queste categorie, sfruttando motori di Intelligenza Artificiale o espressioni regolari.

\subsection{Tassonomia dei tag}
OpenMetadata organizza i tag all'interno di categorie logiche centralizzate, denominate \emph{Classifications}. Per garantire una corretta gestione dei 
metadati e tradurre i requisiti normativi in regole tecniche, la piattaforma fornisce alcune classificazioni di sistema predefinite. Tali categorie non possono essere 
eliminate, ma guidano l'assegnazione corretta delle etichette:
\begin{itemize}
    \item \textbf{PII}: Classificazione standard concepita per l'identificazione e la mappatura delle informazioni personali, fondamentale per la conformità al GDPR. Include il tag \texttt{PII.NonSensitive}, utilizzato per dati identificativi generici e a basso rischio (come nome, cognome o e-mail aziendale), e il tag \texttt{PII.Sensitive}, obbligatorio per gli identificatori critici (come codici fiscali, password o coordinate bancarie). L'assegnazione di quest'ultimo innesca automaticamente logiche di anonimizzazione per gli utenti sprovvisti dei privilegi giusti.
    \item \textbf{PersonalData}: Classificazione complementare per mappare i dati riferibili a persone fisiche, suddivisa in \texttt{PersonalData.Personal} e \texttt{PersonalData.SpecialCategory}. Quest'ultimo tag è destinato ai dati particolari (Art. 9 del GDPR, come dati biometrici o sanitari) e ai dati giudiziari, imponendo un approccio di massima chiusura che ne limita l'accesso esclusivamente alle autorità o ai ruoli aziendali espressamente autorizzati.
    \item \textbf{\emph{Tier} (Livello)}: Utilizzata per suddividere gli asset in base alla loro rilevanza strategica e architetturale per l'organizzazione. Comprende cinque livelli gerarchici standard, da \texttt{Tier.Tier1} (assegnato ad asset contenenti dati critici e insostituibili per il business) fino a \texttt{Tier.Tier5} (associato a informazioni marginali o secondarie).
    \item \textbf{\emph{Certification} (Certificazione)}: Impiegata all'interno delle procedure di Data Quality per qualificare visivamente l'affidabilità di un dataset. Questa tassonomia si articola nei tag gerarchici \texttt{Bronze}, \texttt{Silver} e \texttt{Gold}, permettendo agli analisti di Business Intelligence di verificare all'istante se un dato è stato validato e approvato prima di importarlo nei report.
\end{itemize}

\begin{figure}[H]
    \centering
    \begin{subfigure}[b]{0.9\textwidth}
        \centering
        \includegraphics[width=\textwidth]{schermate/tag_pii.png}
        \caption{Tassonomia PII}
    \end{subfigure}
    
    \vspace{0.4cm}
    
    \begin{subfigure}[b]{0.9\textwidth}
        \centering
        \includegraphics[width=\textwidth]{schermate/tag_tier.png}
        \caption{Tassonomia Tier}
    \end{subfigure}
    
    \caption{Pannelli di gestione delle classificazioni di sistema}
    \label{fig:tassonomie_sistema}
\end{figure}

\noindent Oltre alle tassonomie predefinite, la piattaforma supporta la creazione di classificazioni personalizzate, adattabili alle specifiche policy interne 
di Wintech, per esempio per contrassegnare dati con un tag che ne indica il livello di «confidenzialità».

Ogni classificazione può essere configurata attivando l'opzione di mutua esclusione. Se abilitata, questa proprietà garantisce l'integrità logica della classificazione, 
impedendo che a un singolo asset vengano associati contemporaneamente due tag contrastanti appartenenti alla medesima tassonomia.

\begin{figure}[H]
    \centering
    \includegraphics[width=0.75\textwidth]{schermate/tag_personalizzato.png}
    \caption[Esempio di classificazione personalizzata]{Esempio di classificazione personalizzata \texttt{DataClassification}, configurata con l'opzione di mutua esclusione (\emph{Mutually Exclusive}) attivata}
    \label{fig:tag_personalizzato}
\end{figure}

\subsection{Classificazione automatica}
L'\emph{Auto-Classification} è un workflow di OpenMetadata che rileva e classifica automaticamente i dati sensibili nelle colonne dei database, eliminando la necessità 
di tagging manuale. Questa procedura, eseguita in background e integrabile all'interno delle \emph{pipeline} di sincronizzazione programmata, sfrutta l'operato 
dell'\emph{AutoClassification Agent} per scansionare le sorgenti collegate, applicando o suggerendo etichette come \texttt{PII.Sensitive} e \texttt{PII.NonSensitive}.

Il motore di \emph{auto-classificazione} utilizza due approcci complementari: 
\begin{enumerate}
    \item \textbf{\emph{Column Name Scanner}}: Esegue una scansione strutturale rapida dei metadati, confrontando i nomi fisici delle colonne del database con un elenco di espressioni regolari (regex) predefinite o parole chiave note (identificando, ad esempio, colonne nominate \texttt{email}, \texttt{phone} o \texttt{fiscal\_code}).
    \item \textbf{\emph{Named Entity Recognition} (NER)}: Qualora sia stata abilitata la raccolta di dati campione (\emph{Sample Data}) durante la fase di \emph{ingestion}, il motore NLP analizza il contenuto testuale reale dei record estratti per dedurne semanticamente la natura, identificando le informazioni personali anche in presenza di colonne con nomenclatura generica o non strutturata.
\end{enumerate}

\begin{figure}[H]
    \centering
    \includegraphics[width=1\textwidth]{schermate/dati_sensibili.png}
    \caption[Tabella con tag PII applicati automaticamente]{Esempio di una tabella in cui il motore di auto-classificazione ha applicato automaticamente il tag \texttt{PII.Sensitive} alle colonne contenenti dati personali}
    \label{fig:dati_sensibili}
\end{figure}

\subsubsection{Funzionamento dell'algoritmo di NER}
Il modulo di NER si basa sul \emph{framework} \emph{open-source} \textbf{Microsoft Presidio}, uno strumento specializzato nella rilevazione e anonimizzazione 
dei dati personali che si appoggia a librerie NLP (come spaCy o Stanza). 

In OpenMetadata, Presidio opera come pacchetto Python eseguito direttamente all'interno dell'ambiente di Data Ingestion, gestendo il flusso dei dati in modo \emph{stateless}. 
I dati campione prelevati dal database vengono caricati temporaneamente nella RAM del Container ed elaborati dal modello, il quale restituisce un file in formato JSON 
contenente le informazioni rilevate. Al termine della scansione, il campione viene immediatamente distrutto dalla memoria e l'agente trasmette al server centrale 
esclusivamente il metadato risultante (il tag di classificazione da applicare).

\begin{figure}[H]
    \centering
    \includegraphics[width=0.75\textwidth]{autoclassificazione.png}
    \vspace{0.2cm}
    \caption[Rappresentazione del flusso di auto-classificazione]{Architettura dell'auto-classificazione: i dati reali vengono elaborati e distrutti localmente, trasmettendo al server centrale solo i metadati. \textit{Nota: La presente infografica è stata generata con l'ausilio di un sistema di Intelligenza Artificiale}}
    \label{fig:flusso_autoclassificazione}
\end{figure}

Per rispettare il GDPR e i vincoli di sicurezza imposti dall'infrastruttura di Wintech, l'intero processo avviene in locale e offline. A differenza degli assistenti basati 
su Cloud, i modelli linguistici necessari per l'analisi vengono scaricati una sola volta durante la fase iniziale di compilazione e installazione del Container. 
A riprova di questo, su reti aziendali dotate di rigide policy di sicurezza, il Container può restituire l'errore \texttt{SSL: CERTIFICATE\_VERIFY\_FAILED} 
al primo avvio, venendo bloccato dal firewall aziendale durante il tentativo di download del modello; una volta scaricato manualmente, il sistema opera totalmente 
offline garantendo che nessuna informazione sensibile venga trasmessa all'esterno della rete aziendale.

Infine, l'architettura basata su Presidio garantisce una notevole scalabilità tramite i cosiddetti \emph{Custom Recognizers}. Qualora l'azienda necessiti di intercettare 
formati di dati specifici o interni (per esempio la sintassi proprietaria di un ID Cliente), è sufficiente sviluppare e integrare un riconoscitore personalizzato 
all'interno della \emph{pipeline} Python, estendendo le capacità di rilevamento senza alterare i protocolli di sicurezza.


\subsection{Limiti dell'automazione e classificazione manuale}
Durante i test di implementazione, è emerso che la funzionalità di classificazione automatica non è esente da imprecisioni contestuali e presenta alcune limitazioni architetturali. 

In primo luogo, l'efficacia del rilevamento automatico sui database aziendali dipende dall'ottimizzazione di parametri tecnici specifici:
\begin{itemize}
    \item \textbf{Propagazione dei prerequisiti logici}: l'esecuzione dell'\emph{AutoClassification Agent} è strettamente subordinata all'avvenuto completamento delle attività del \emph{Profiler Agent}. Poiché il motore NLP necessita del contenuto testuale reale dei record (\emph{Sample Data}) per applicare la Named Entity Recognition, è indispensabile pianificare l'esecuzione delle \emph{pipeline} garantendo che il profilo dei dati venga generato prima della scansione di classificazione;
    \item \textbf{Regolazione del livello di fiducia (\emph{confidence level})}: per bilanciare la precisione del rilevamento, il sistema consente di calibrare una soglia minima di confidenza. Questo parametro è fondamentale per evitare sia l'insorgere di falsi positivi (dati innocui scambiati per sensibili) sia di falsi negativi (dati sensibili non rilevati).
\end{itemize}

In secondo luogo, è stato rilevato un importante vincolo relativo alla versione \emph{open-source} di OpenMetadata: il modulo è progettato per operare una classificazione 
prevalentemente binaria a livello di GDPR, etichettando i campi sensibili in modo generico come \texttt{PII.Sensitive} o \texttt{PII.NonSensitive}. La capacità di associare 
automaticamente tag più granulari (come \texttt{Address} o \texttt{Name}) è una funzionalità esclusiva della distribuzione commerciale (\emph{Collate}) e non è direttamente 
fruibile nell'installazione locale standard. 

A fronte di questi limiti tecnici e per garantire la massima conformità alle policy aziendali, è indispensabile una fase di classificazione e revisione manuale. 
In questo passaggio, figure qualificate come i \emph{data steward} o i proprietari del dato (\emph{data owners}) intervengono per validare, correggere o raffinare le 
etichette suggerite dalla piattaforma. Questa attività permette anche di arricchire i metadati dell'asset: l'utente può inserire descrizioni contestuali dettagliate, 
definire l'appartenenza a specifici domini aziendali e associare i campi ai termini formalizzati nel glossario.

\subsection{Integrazione dell'AI Generativa tramite MCP}
A completamento delle analisi architetturali svolte durante lo stage, si è esplorata la possibilità di integrare strumenti di Intelligenza Artificiale generativa 
in Cloud (come ChatGPT o Claude) per abilitare la ricerca semantica e l'interrogazione del catalogo in linguaggio naturale. Questo è possibile grazie all'integrazione 
nativa all'interno di OpenMetadata, disponibile anche nella versione \emph{open-source}, di un server basato sullo standard \textbf{\gls{mcp}}. 

Il protocollo MCP fa da interfaccia di comunicazione sicura tra il catalogo dati e gli assistenti AI esterni. Condividendo esclusivamente i metadati, 
questa soluzione risponde ai requisiti di sicurezza attraverso i seguenti meccanismi:

\begin{itemize}
    \item \textbf{Totale isolamento dei dati reali}: L'assistente AI non stabilisce alcuna connessione diretta con i database fisici dell'azienda, precludendo 
    l'accesso ai dati in chiaro. Il modello interagisce unicamente con il server MCP di OpenMetadata, il quale trasmette all'esterno esclusivamente metadati 
    strutturali (schema delle tabelle, nomi dei campi, definizioni del glossario, tag di classificazione applicati). Le informazioni effettive contenute nei 
    database non vengono mai elaborate dall'AI e non lasciano il perimetro aziendale, garantendo la conformità alle direttive sulla sicurezza.
    \item \textbf{Preservazione del controllo degli accessi}: L'utilizzo dell'AI tramite server MCP è subordinato al rispetto delle policy di Data Governance, impedendo 
    all'Intelligenza Artificiale di bypassare i sistemi di sicurezza interni. La connessione al server richiede un'autenticazione sicura: l'assistente AI eredita 
    direttamente i permessi dell'utente che ha effettuato l'accesso, seguendo il modello di controllo degli accessi basato sui ruoli (RBAC) configurato sulla piattaforma. 
    Se l'account dell'utente che interroga l'AI non dispone delle autorizzazioni necessarie per visualizzare una determinata area del catalogo, il server MCP bloccherà 
    la richiesta alla fonte, impedendo all'AI di accedere ai relativi metadati e di includerli nella risposta.
\end{itemize}


\section{Mascheramento dei dati}
A differenza di quanto ipotizzato in fase di analisi dei requisiti, la piattaforma non dispone di impostazioni personalizzabili di anonimizzazione o pseudonimizzazione dei dati. 
Tuttavia, OpenMetadata integra un meccanismo automatico di mascheramento totale (anonimizzazione) strettamente dipendente dal modello di accesso basato sui ruoli.

Quando una colonna o una tabella viene classificata con il tag \texttt{PII.Sensitive}, i relativi \emph{Sample Data} vengono oscurati e sostituiti visivamente da asterischi 
(\texttt{******}) all'interno dell'interfaccia. Questo blocco visivo viene applicato in automatico a tutti gli utenti standard, garantendo invece la visualizzazione dei valori 
in chiaro esclusivamente agli amministratori di sistema (\emph{Admin}) e ai proprietari (\emph{owner}) del singolo asset.

\begin{figure}[H]
    \centering
    \includegraphics[width=1\textwidth]{schermate/anonimizzazione.png}
    \vspace{0.2cm}
    \caption[Applicazione dell'anonimizzazione nell'interfaccia]{Applicazione pratica dell'anonimizzazione: i \emph{Sample Data} classificati come sensibili vengono oscurati per gli utenti non autorizzati}
    \label{fig:anonimizzazione}
\end{figure}

\subsection{Limiti architetturali}
Durante l'analisi e il collaudo della funzionalità, sono state riscontrate alcune limitazioni:
\begin{itemize}
    \item \textbf{Assenza di pseudonimizzazione}: Non è possibile scegliere tra anonimizzazione e pseudonimizzazione. Il comportamento del meccanismo di mascheramento è predefinito e si applica in modo uniforme (mascheramento totale) a tutti gli asset etichettati come \texttt{PII.Sensitive}.
    \item \textbf{Rigidità configurativa}: Non è permesso configurare autonomamente regole di mascheramento basate sull'assegnazione di tag di classificazione personalizzati.
    \item \textbf{Mancata propagazione esterna}: Non sono previste funzionalità per propagare le direttive di mascheramento ai sistemi esterni a valle (come, ad esempio, gli strumenti di Business Intelligence).
    \item \textbf{Dipendenza dal modello RBAC}: Il meccanismo si basa interamente sul sistema di \emph{ownership} e ruoli. Una rigorosa configurazione dei permessi rappresenta quindi un prerequisito architetturale indispensabile affinché la protezione dei dati sensibili sia garantita.
\end{itemize}

Nonostante questi vincoli tecnici, OpenMetadata offre una funzionalità di base per l'anonimizzazione dei dati che risulta comunque sufficiente per soddisfare i requisiti normativi e le esigenze aziendali.


\section{Gestione degli accessi}
L'efficacia della Data Governance e la conformità al GDPR richiedono una gestione degli accessi sicura, granulare e centralizzata. L'adozione di un Data Catalog 
aziendale deve garantire il rispetto del \emph{Principle of Least Privilege}, fungendo da unico punto di controllo per l'accesso ai metadati e proteggendo in particolare le informazioni sensibili.

Per rispondere a questa esigenza, OpenMetadata utilizza un modello di controllo degli accessi basato sui ruoli (RBAC), combinato con una fedele mappatura 
gerarchica dell'organigramma aziendale, possibile grazie alla definizione di utenti e team.

\subsection{Mappatura dell'organigramma: utenti e team}
Il primo livello di sicurezza riguarda la gestione delle identità. Dalla scheda di amministrazione è possibile configurare i singoli utenti (\emph{User}) e 
raggrupparli in gruppi di lavoro (\emph{Team}). Questa struttura permette di replicare fedelmente l'organigramma di Wintech all'interno del sistema. 
OpenMetadata permette di organizzare l'albero aziendale secondo cinque livelli logici (\emph{teamType}):
\begin{itemize}
    \item \textbf{\emph{Organization}}: È il nodo radice della gerarchia. Contiene tutti i dipendenti e i sottogruppi dell'azienda e non possiede nodi genitori.
    \item \textbf{\emph{BusinessUnit}}: Rappresenta le divisioni di primo livello dell'azienda (es. Direzione, Ricerca e Sviluppo). Può dipendere esclusivamente dall'\emph{Organization} o da un'altra \emph{BusinessUnit}.
    \item \textbf{\emph{Division}}: È una partizione operativa intermedia. Dipende gerarchicamente da un solo nodo superiore (\emph{Organization}, \emph{BusinessUnit} o un'altra \emph{Division}).
    \item \textbf{\emph{Department}}: Individua aree di competenza specialistiche. A differenza degli altri livelli, può ereditare proprietà da più nodi genitori, offrendo maggiore flessibilità organizzativa.
    \item \textbf{\emph{Group}}: Rappresenta il nodo foglia della gerarchia. Non può contenere sottogruppi, ma è formato solo da singoli \emph{User}. Nel Data Catalog, solo i team di tipo \emph{Group} possono detenere formalmente la proprietà (\emph{Ownership}) degli asset.
\end{itemize}

L'inserimento degli utenti può avvenire manualmente, ma l'infrastruttura supporta anche la centralizzazione degli accessi tramite l'integrazione con protocolli di 
Single Sign-On (SSO), forniti da provider come Microsoft Entra ID, Google Workspace e Okta. 
Questa configurazione abilita il provisioning automatico (\emph{Auto-Provisioning}). Al momento del primo accesso di un utente tramite credenziali aziendali, 
la piattaforma interroga il provider e crea in automatico il relativo profilo nel catalogo, riducendo l'onere amministrativo per la gestione del personale.

\begin{figure}[H]
    \centering
    \begin{subfigure}[b]{1\textwidth}
        \centering
        \includegraphics[width=\textwidth]{schermate/user.png}
        \caption{Pannello di amministrazione dei singoli utenti (\emph{User})}
    \end{subfigure}
    
    \vspace{0.4cm}
    
    \begin{subfigure}[b]{1\textwidth}
        \centering
        \includegraphics[width=\textwidth]{schermate/team.png}
        \caption{Mappatura gerarchica dell'organigramma tramite i \emph{Team}}
    \end{subfigure}
    
    \caption[Pannelli di gestione degli utenti e dei team]{Interfacce per la gestione delle identità in OpenMetadata: visualizzazione dell'elenco utenti (sopra) e dell'organizzazione gerarchica in dipartimenti o \emph{BusinessUnit} (sotto)}
    \label{fig:gestione_utenti_team}
\end{figure}

\subsection{L'architettura Role-Based Access Control}
Una volta definita la struttura degli utenti e dei team, la piattaforma permette di definire autorizzazioni e permessi attraverso il sistema \gls{rbac}.
Dal punto di vista architetturale, il motore di controllo degli accessi si basa su tre componenti interconnessi:
\begin{enumerate}
    \item \textbf{Rules (Regole)}: Sono l'elemento base della sicurezza. Definiscono se una specifica azione (es. modificare tag, visualizzare dati campione) è permessa o negata su un certo tipo di risorsa.
    \item \textbf{Policies (Politiche)}: Sono raggruppamenti logici di regole. Vengono create per rispondere a specifiche esigenze di sicurezza aziendale (es. policy per la qualità dei dati, policy per la gestione dei domini).
    \item \textbf{Roles (Ruoli)}: Associano una o più politiche ai profili utente. Assegnando un ruolo a un utente o a un intero team, le relative autorizzazioni vengono propagate istantaneamente a tutti i membri.
\end{enumerate}

\subsection{Logica di valutazione dei permessi}
Quando un utente tenta di eseguire un'operazione su un asset, il motore di OpenMetadata valuta la richiesta. La decisione di autorizzare o bloccare l'azione segue 
quattro principi fondamentali:
\begin{itemize}
    \item \textbf{Permessi di base (Default)}: Ogni utente registrato ottiene i privilegi minimi associati al ruolo base (\emph{Data Consumer Role}) e alla policy globale (\emph{OrganizationPolicy}). Questo garantisce che chiunque possa cercare e leggere i dati (Data Discovery), escludendo però le operazioni di modifica.
    \item \textbf{Ereditarietà dai team}: I permessi sono cumulativi. Il sistema calcola l'unione logica tra i privilegi assegnati direttamente al profilo e quelli ereditati da tutti i team a cui l'utente appartiene.
    \item \textbf{Eccezione di Proprietà (\emph{Ownership})}: È una deroga fondamentale alle regole standard. Se un utente o un team è designato come proprietario di un asset, ne ottiene automaticamente la gestione completa, superando i limiti del proprio ruolo ordinario.
    \item \textbf{Priorità dei divieti (\emph{Deny over Allow})}: È la regola di sicurezza che previene i conflitti. Le regole all'interno delle policy possono consentire (\emph{Allow}) o vietare (\emph{Deny}) un'azione. Se un utente eredita un permesso abilitante, ma appartiene a un team con una regola restrittiva esplicita (\emph{Deny}) per la stessa azione, il divieto prevale sempre e l'operazione viene bloccata.
\end{itemize}

\begin{figure}[H]
    \centering
    \begin{subfigure}[b]{1\textwidth}
        \centering
        
\includegraphics[width=\textwidth]{schermate/ruolo_datagovmanager.png}
        \caption{Assegnazione di una \emph{policy} di sicurezza allo specifico ruolo di \emph{Data Governance Manager}}
    \end{subfigure}
    
    \vspace{0.4cm}
    
    \begin{subfigure}[b]{1\textwidth}
        \centering
        
\includegraphics[width=\textwidth]{schermate/policy_gdpr.png}
        \caption{Dettaglio della \emph{policy}: regola con effetto di divieto (\emph{Deny}) basata sui tag}
    \end{subfigure}
    
    \caption[Esempio di architettura RBAC per la conformità GDPR]{Applicazione pratica dell'architettura RBAC: il Ruolo (sopra) eredita la \emph{policy} dedicata al rispetto del GDPR, la quale (sotto) definisce una regola di divieto (\emph{Deny}) per impedire la visualizzazione in chiaro dei dati etichettati come sensibili}
    \label{fig:rbac_gdpr}
\end{figure}


\section{Glossario aziendale}
Un \textbf{Glossario aziendale} (\emph{Business Glossary}) è un vocabolario centralizzato che raccoglie, definisce e standardizza i concetti e la terminologia 
rilevante all'interno dell'organizzazione. Rappresenta uno strumento di Data Governance fondamentale, a cui ogni utente può fare riferimento per risolvere ambiguità 
semantiche e comprendere appieno il significato del patrimonio informativo.

L'obiettivo principale di questo strumento è stabilire un linguaggio aziendale comune, colmando la distanza comunicativa tra i diversi reparti 
attraverso la definizione di un significato comune per ogni termine. All'interno di OpenMetadata, i concetti formalizzati nel glossario assumono il ruolo di veri e propri 
metadati: possono essere infatti assegnati agli asset per descrivere e categorizzare ulteriormente tabelle e dataset.

Un'organizzazione può definire molteplici glossari indipendenti, strutturandoli in base agli specifici domini aziendali (ad esempio Amministrazione, Risorse Umane o Vendite). 
I termini all'interno di ciascun glossario possono essere organizzati secondo una gerarchia logica e possono essere arricchiti con una serie di proprietà avanzate:
\begin{itemize}
    \item \textbf{Contesto semantico}: descrizioni testuali, sinonimi, riferimenti a link esterni e associazioni con altri termini correlati;
    \item \textbf{Propagazione delle classificazioni}: possibilità di associare tag di classificazione direttamente al termine. Quando il concetto viene applicato a un dataset, quest'ultimo eredita automaticamente anche i tag associati al termine stesso;
    \item \textbf{Vincoli di assegnazione}: funzionalità di mutua esclusione (\emph{Mutually Exclusive}), che impedisce di assegnare simultaneamente due termini concettualmente in conflitto al medesimo asset;
    \item \textbf{Validazione e ownership}: designazione di specifici utenti revisori (\emph{reviewer}), incaricati di validare i termini, approvare le modifiche e garantire l'accuratezza del glossario nel tempo.
\end{itemize}

\begin{figure}[H]
    \centering
    \includegraphics[width=1\textwidth]{schermate/glossario_wintech.png}
    \vspace{0.2cm}
    \caption{Interfaccia principale del Glossario aziendale}
    \label{fig:glossario_generale}
\end{figure}

\begin{figure}[H]
    \centering
    \includegraphics[width=1\textwidth]{schermate/glossario_termini.png}
    \vspace{0.2cm}
    \caption{Struttura gerarchica dei termini nel Glossario}
    \label{fig:glossario_gerarchia}
\end{figure}

\section{Domini e data products}
OpenMetadata permette di organizzare gli asset di dati in gruppi logici, definiti Domini (\emph{domain}), per strutturare il patrimonio informativo in modo più efficace. 
Un Dominio consente di categorizzare e gestire asset, glossari, team e altre entità sotto una struttura logica unificata, supportando un'architettura decentralizzata in 
cui ogni area è gestita in modo indipendente con un'\emph{ownership} distribuita.

\subsection{Valore aziendale e Data Product}
L'introduzione dei Domini permette di decentralizzare la gestione dei dati: delegando la responsabilità ai singoli dipartimenti (come Risorse Umane o Marketing), 
la cura e la documentazione del dato vengono affidate a chi ne possiede le reali competenze. Questo approccio elimina i colli di bottiglia e agevola l'applicazione 
delle policy di sicurezza.

All'interno di ciascun Dominio, gli asset possono essere ulteriormente raggruppati in \textbf{Data Product}. Un Data Product è semplicemente un sottoinsieme di 
dati organizzato per rispondere a un'esigenza specifica (ad esempio, raggruppare solo i dati relativi alle campagne pubblicitarie). Questa suddivisione trasforma 
i dati grezzi in un vero e proprio «prodotto» pronto all'uso, facilmente consultabile dagli utenti o analizzabile tramite strumenti di Business Intelligence.

\begin{figure}[H]
    \centering
    \begin{subfigure}[b]{1\textwidth}
        \centering
        \includegraphics[width=\textwidth]{schermate/domini.png}
        \caption{Vista generale dei Domini configurati nella piattaforma}
    \end{subfigure}
    
    \vspace{0.4cm}
    
    \begin{subfigure}[b]{1\textwidth}
        \centering
        \includegraphics[width=\textwidth]{schermate/data_products.png}
        \caption{Dettaglio dei Data Product (\emph{Advertising} e \emph{Social Media Management}) nel dominio Marketing}
    \end{subfigure}
    
    \caption{Organizzazione logica in Domini e Data Product}
    \label{fig:domini_data_products}
\end{figure} %7
    % capitolo 8: integrazione della business intelligence
\chapter{Integrazione con la Business Intelligence}
\label{cap:business-intelligence}

\intro{Questo capitolo mette in relazione i due temi studiati durante il percorso di stage, il catalogo dati e la Business Intelligence, delineando il flusso ideale e sicuro delle informazioni.}\\

\section{Introduzione alla Business Intelligence}
La Business Intelligence (BI) è una disciplina che combina infrastrutture tecnologiche, strumenti di analisi dei dati, visualizzazioni avanzate e \emph{best practice} 
organizzative con l'obiettivo di supportare le aziende nell'adozione di decisioni basate su evidenze empiriche ed oggettive (\emph{data-driven decision making})\cite{site:business-intelligence}.

Dal punto di vista teorico e operativo, il flusso della Business Intelligence si articola in quattro passaggi sequenziali fondamentali:
\begin{enumerate}
    \item \textbf{Raccolta ed estrazione}: i dati grezzi provenienti da fonti eterogenee (quali database, fogli di calcolo Excel o sistemi \emph{on-premise}) vengono estratti e centralizzati in depositi unificati;
    \item \textbf{Analisi descrittiva}: le informazioni memorizzate vengono elaborate per comprendere l'andamento storico e corrente dell'azienda. In questa fase si opera una distinzione essenziale rispetto alla Business Analytics (BA): mentre la Business Intelligence si concentra sull'analisi descrittiva (rispondendo alle domande \emph{cosa} e \emph{come}), la Business Analytics dà priorità all'analisi predittiva e prescrittiva, sfruttando modelli statistici per prevedere eventi futuri (rispondendo alla domanda \emph{perché});
    \item \textbf{Visualizzazione dei dati}: i risultati dell'elaborazione analitica vengono convertiti in rappresentazioni grafiche intuitive (istogrammi, diagrammi a torta, \emph{scatter plot}), rendendo le metriche complesse accessibili e leggibili anche per gli utenti di business non tecnici;
    \item \textbf{Azione basata su informazioni in tempo reale}: la visualizzazione delle performance storiche e attuali, calata nel contesto aziendale, consente ai \emph{decision maker} di intervenire prontamente per ottimizzare i processi operativi.
\end{enumerate}

\begin{figure}[H]
    \centering
    \includegraphics[width=1\textwidth]{bi/bi_concept.png}
    \caption[Flusso logico della Business Intelligence]{Rappresentazione del flusso logico della Business Intelligence. \textit{Nota: La presente infografica è stata generata con l'ausilio di un sistema di Intelligenza Artificiale}}
    \label{fig:bi_concept}
\end{figure}

All'interno dell'infrastruttura di Wintech, la piattaforma di riferimento per l'analisi aziendale e la Business Intelligence è Microsoft Power BI\cite{site:powerbi-docs}. 
Questa soluzione fornisce servizi e strumenti in grado di connettere, trasformare e condividere i dati dell'organizzazione, traducendoli in informazioni dettagliate e visualizzazioni grafiche.

Utilizzato quotidianamente dall'azienda, Power BI soddisfa nativamente numerosi requisiti legati all'osservabilità e alla gestione del dato, offrendo:
\begin{itemize}
    \item \textbf{Centralizzazione delle fonti}: la capacità di connettere in un unico ambiente un'ampia varietà di dati strutturati (come database SQL Server, PostgreSQL) e non strutturati (come file JSON, CSV, fogli Excel), fungendo da singola fonte di verità per la reportistica aziendale;
    \item \textbf{Sincronizzazione programmata}: l'allineamento costante e automatico delle informazioni con le fonti originali, tramite Power BI Service e i \emph{Gateway} dati locali;
    \item \textbf{Organizzazione logica}: la strutturazione dei modelli di dati attraverso prospettive distinte (vista Report, vista Tabella, vista Modello), che permettono di esplorare le relazioni tra le entità importate;
    \item \textbf{Sicurezza e gestione degli accessi}: l'implementazione di standard di sicurezza, includendo un sistema RBAC a livello di \emph{workspace} e la sicurezza a livello di riga (RLS - \emph{Row-Level Security}).
\end{itemize}

\subsection{La piattaforma Microsoft Power BI}
Wintech adotta Microsoft Power BI come soluzione di riferimento per lo sviluppo dei propri sistemi di reporting interni e per la consulenza offerta ai clienti. 
Si tratta di una piattaforma avanzata di Business Intelligence che fornisce una \emph{suite} di servizi, applicazioni e connettori che trasformano molteplici 
fonti di dati in informazioni strategiche, dettagliate e visivamente interattive.

La piattaforma offre tre ambienti operativi principali, che supportano l'analista durante l'intero ciclo di modellazione e fruizione dei dati:
\begin{itemize}
    \item \textbf{Power BI Desktop}: Applicazione client di sviluppo utilizzata per connettersi alle sorgenti, effettuare la pulizia e la modellazione dei dati tramite il Power Query Editor, definire le misure in linguaggio DAX (\emph{Data Analysis Expressions}) e creare i report visivi.
    \item \textbf{Power BI Service}: Servizio in cloud basato su modello SaaS (\emph{Software as a Service}) che permette la pubblicazione, condivisione e distribuzione sicura dei report all'interno dell'organizzazione, nonché la creazione di dashboard riepilogative destinate al \emph{management}.
    \item \textbf{Power BI Mobile}: Applicazioni per dispositivi mobili ottimizzate per visualizzare i report in mobilità, sfruttando specifici layout reattivi.
\end{itemize}

\begin{figure}[H]
    \centering
    
\includegraphics[width=0.25\textwidth]{logo-powerbi.png}
    \caption{Logo ufficiale di Microsoft Power BI}
    \label{fig:power_bi_logo}
\end{figure}

\subsection{Contesto aziendale}
In Wintech S.p.A., la Business Intelligence ha un ruolo strategico. In qualità di consulente e \emph{System Integrator}, l'azienda aiuta i clienti a integrare 
sistemi ERP e a creare dashboard personalizzate, che permettono di monitorare facilmente le vendite, la logistica, le risorse umane e l'andamento finanziario.

Durante lo stage, lo studio di Power BI ha permesso di comprendere l'utilità pratica dello strumento, sviluppando report basati sia sul database di test \emph{AdventureWorks} 
sia sui movimenti contabili reali estratti in ambiente sicuro dai sistemi gestionali aziendali.

A titolo esemplificativo, la Figura~\ref{fig:report_powerbi} mostra l'interfaccia di un report sviluppato sui dati di test, evidenziando le funzionalità di 
Data Visualization e il pannello per la configurazione degli avvisi di monitoraggio.

\begin{figure}[H]
    \centering
    \includegraphics[width=1\textwidth]{bi/report_powerbi.png}
    \caption[Interfaccia di un report in Microsoft Power BI]{Esempio di dashboard analitica sviluppata in Power BI sui dati di \emph{AdventureWorks}. Sulla destra è visibile il pannello di configurazione degli avvisi, utile per il monitoraggio parziale dei dati}
    \label{fig:report_powerbi}
\end{figure}

Tuttavia, Power BI è uno strumento progettato specificamente per la Data Visualization e la reportistica. Proprio per questo motivo, elabora le informazioni 
esattamente come vengono estratte, senza verificarne preventivamente la qualità o la presenza di dati personali alla sorgente. 

Affinché la Business Intelligence sia davvero affidabile e conforme alle normative sulla privacy, è necessario creare un flusso informativo intermedio. 
Questo livello si posiziona tra le fonti grezze e i report finali, affidando l'ispezione, la pulizia e la classificazione dei dati a piattaforme specifiche di Data Discovery e Data Governance.


\section{Flusso sicuro dei dati}
% PER SPIEGARE QUAL è L'IDEA DIETRO LA CONNESSIONE DATA CATALOG - POWER BI + grafico del flusso:  dati grezzi → catalogazione → classificazione PII → data quality e 
% monitoraggio → mascheramento → governance e policy → strumenti BI che consumano solo dati puliti, sicuri e certificati
L'efficacia strategica della Business Intelligence non dipende solo dalle capacità analitiche o grafiche delle dashboard, ma è condizionata alla radice dall'accuratezza e 
dalla sicurezza del flusso informativo sottostante. Nelle architetture tradizionali, il collegamento diretto tra gli strumenti di Business Intelligence (come Microsoft 
Power BI) e le basi dati genera criticità significative sotto il profilo operativo, della Data Quality e della conformità al GDPR.

Connettere gli strumenti di Business Intelligence direttamente alle sorgenti espone l'organizzazione a tre vulnerabilità principali:
\begin{itemize}
    \item \textbf{Esposizione incontrollata di dati sensibili}: i report di BI analizzano tipicamente metriche aggregate. Tuttavia, una connessione diretta ai database di produzione consente agli analisti di accedere e importare record contenenti informazioni personali in chiaro (quali nomi, indirizzi, e-mail o codici fiscali), violando i principi di minimizzazione e limitazione delle finalità previsti dal GDPR;
    \item \textbf{Propagazione di anomalie di Data Quality}: se una sorgente dati presenta valori nulli, duplicati o anomalie strutturali, lo strumento di BI importa passivamente tali errori. Questo genera visualizzazioni distorte o errate, compromettendo l'affidabilità delle metriche esposte e quindi dei processi decisionali;
    \item \textbf{Assenza di contesto semantico e documentale}: gli analisti faticano a individuare e comprendere i dataset d'interesse (Data Discovery), dovendo decifrare strutture fisiche complesse e nomenclature tecniche prive di una descrizione formalizzata o di un'associazione chiara con i termini di business dell'organizzazione.
\end{itemize}

Per mitigare tali rischi e implementare nativamente i principi di \emph{Privacy by Design}, l'architettura informativa deve integrare un livello intermedio che funga 
da \emph{gatekeeper} logico. Un Data Catalog come OpenMetadata si interpone tra i sistemi gestionali e le piattaforme analitiche, orchestrando il ciclo di 
vita dei metadati e assicurando che la Business Intelligence consumi esclusivamente risorse certificate e sicure.

Il processo che trasforma i dati grezzi in informazioni utili per le decisioni aziendali si articola in sei fasi:
\begin{enumerate}
    \item \textbf{Estrazione dei metadati}: il modulo di Data Ingestion del catalogo interroga periodicamente le sorgenti in modalità \emph{Read-Only}, estraendo la sola topologia strutturale delle tabelle (nomi, colonne, tipi di dato) e un campione ristretto di record di esempio (\emph{Sample Data}), senza duplicare fisicamente i dati reali o appesantire l'infrastruttura transazionale;
    \item \textbf{Arricchimento semantico}: i metadati tecnici vengono associati ai termini formalizzati nel glossario aziendale, definendo responsabilità e \emph{ownership} per ciascun asset e facilitando la Data Discovery da parte degli utenti di business;
    \item \textbf{Individuazione e classificazione dei PII}: mediante l'esecuzione dell'\emph{AutoClassification Agent} o l'applicazione manuale dei tag, il sistema rileva le informazioni personali e vi applica i tag pertinenti (es. \texttt{PII.Sensitive}), mappando la distribuzione dei dati protetti;
    \item \textbf{Monitoraggio della qualità}: il motore di Data Quality esegue i test pianificati sui database sorgente. In caso di anomalie, il sistema segnala l'errore sulla dashboard di Data Observability e notifica i responsabili tramite alert automatici, impedendo la propagazione dell'errore a valle;
    \item \textbf{Mascheramento e anonimizzazione}: nel catalogo l'accesso ai dati sensibili del \emph{Sample Data} viene protetto tramite le policy RBAC, oscurando i valori (\texttt{******}) ai profili non autorizzati. A livello di database fisico, i dataset destinati ad ambienti di test e reporting vengono opportunamente anonimizzati tramite tecniche e modelli formali, mantenendo i soli vincoli relazionali necessari alle query di analisi;
    \item \textbf{Consumo sicuro}: lo sviluppatore o l'analista di BI si connette esclusivamente alle sorgenti o alle aree di \emph{staging} sicure contenenti i dati precedentemente catalogati, verificati e bonificati, garantendo una reportistica attendibile, accurata e conforme alle normative.
\end{enumerate}

\begin{figure}[H]
    \centering
    
\includegraphics[width=1\textwidth]{bi/flusso_dati.png}
    \caption{Rappresentazione logica del flusso sicuro delle informazioni}
    \label{fig:flusso_dati}
\end{figure}

\subsection{Vincoli di integrazione}
Durante lo stage, a causa di vincoli legati alla disponibilità di licenze Microsoft Power BI Pro e alle policy di rete interne necessarie per l'esposizione delle API, 
non è stato possibile collaudare la connessione fisica tra l'istanza locale di OpenMetadata e Power BI. 
Tali requisiti, di competenza dell'area IT Infrastructure, sono stati posticipati per motivi di sicurezza perimetrale, portando lo stage a concentrarsi sullo studio 
concettuale e sulla modellazione logica del flusso informativo.

Un altro vincolo tecnico consiste nell'asimmetria dell'integrazione tra i due sistemi. All'interno di OpenMetadata è possibile registrare Power BI come un servizio sorgente 
per estrane i metadati, importando nel catalogo i metadati relativi a dashboard, report, datamodel e Data Lineage.

Tuttavia, questo flusso di metadati è esclusivamente unidirezionale:
\begin{itemize}
    \item Power BI alimenta il data catalog fornendo la topologia delle proprie risorse analitiche;
    \item Power BI non può in alcun modo connettersi a OpenMetadata per interrogarlo come fonte dati reale o per estrarre direttamente tabelle fisiche pulite e anonimizzate.
\end{itemize}

Questo limite deriva dall'architettura del catalogo dati: la piattaforma non opera come database fisico e non fa una copia delle informazioni aziendali, ma si limita 
ad archiviare e indicizzare i metadati descrittivi e qualitativi degli asset esterni. Di conseguenza, la pulizia e il mascheramento fisico dei dati devono essere delegati 
ad altri strumenti, che operano a livello di base dati o staging area, lasciando a OpenMetadata il solo ruolo di tracciamento e controllo di conformità.



\section{Strumenti di Data Discovery}
Durante l'analisi condotta nel percorso di stage, un quesito strategico fondamentale ha guidato la valutazione dell'architettura aziendale: è strettamente indispensabile 
introdurre una piattaforma di data catalog complessa e centralizzata come OpenMetadata, o è possibile soddisfare i requisiti di Data Governance combinando le funzionalità 
di Microsoft Power BI con strumenti verticali e leggeri di PII Discovery?

Power BI, ampiamente integrato nel perimetro operativo di Wintech S.p.A., si attesta come un eccellente strumento di reportistica e visualizzazione, ma implementa nativamente 
anche alcune funzionalità tipiche di un catalogo dati:
\begin{itemize}
    \item \textbf{Data Ingestion:} si connette a una vastissima gamma di fonti dati tramite connettori predefiniti;
    \item \textbf{Data Discovery:} permette l'esplorazione dei modelli semantici e la ricerca visiva dei dataset;
    \item \textbf{Data Collaboration:} supporta la condivisione sicura dei report attraverso l'integrazione con Microsoft Teams e Microsoft 365;
    \item \textbf{Data Observability:} consente il monitoraggio dei KPI e l'attivazione di allarmi in tempo reale al variare delle metriche aziendali.
\end{itemize}

Tuttavia, la business intelligence da sola non è in grado di garantire il controllo qualitativo e la conformità normativa alla sorgente. Power BI non effettua test di 
Data Quality automatici sui database sorgente e non dispone di algoritmi per scovare le informazioni personali nei database relazionali aziendali prima che queste vengano 
importate nel modello. 

Per colmare questo divario senza incorrere nella complessità sistemistica di un data catalog completo, lo stage ha previsto la valutazione di cinque soluzioni specialistiche 
(\emph{tool} di PII Discovery e anonimizzazione), analizzandone l'utilità, le caratteristiche e la fattibilità di integrazione nell'infrastruttura aziendale:
\begin{enumerate}
    \item \textbf{Microsoft Presidio}: un \emph{framework} \emph{open-source} e \emph{stateless} sviluppato da Microsoft, specializzato nell'individuazione e nell'anonimizzazione delle PII tramite modelli NLP locali (spaCy) ed espressioni regolari. Poiché non si connette direttamente alle fonti dati ma riceve testo tramite API, risulta estremamente flessibile per essere integrato all'interno di pipeline Python personalizzate prima del caricamento in Power BI;
    \item \textbf{ARX (Data Anonymization Tool)}: uno strumento \emph{open-source} dedicato alla de-identificazione fisica dei dati strutturati. Si distingue per la capacità di connettersi ai database relazionali, analizzare preventivamente il rischio statistico di re-identificazione e applicare modelli matematici formali di privacy, quali la \emph{k-Anonymity}, la \emph{l-Diversity} e la \emph{t-Closeness}, esportando tabelle bonificate e sicure utilizzabili dagli analisti di business intelligence;
    \item \textbf{Hawk-Eye}: una potente utility a riga di comando \emph{open-source} progettata per scansionare sia fonti strutturate (database SQL) sia risorse non strutturate (file PDF, Word, immagini e video) grazie all'uso combinato di pattern matching e tecnologie OCR. Pur eccellendo nella fase di Data Discovery e auditing, Hawk-Eye presenta due limiti significativi: non supporta nativamente Microsoft SQL Server e non applica alcuna tecnica di mascheramento, delegando la protezione a passaggi successivi dello sviluppo;
    \item \textbf{Bytebase}: uno strumento focalizzato sul DevSecOps e sulla gestione del ciclo di vita dei database. Consente l'applicazione di policy di mascheramento dinamico (\emph{dynamic data masking}) direttamente a livello di query, ma richiede l'adozione di un'infrastruttura di gestione complessa;
    \item \textbf{Fides (di Ethyca)}: una piattaforma di privacy engineering basata sulla metodologia \emph{Privacy-as-Code}, progettata per automatizzare il rispetto del GDPR direttamente nei flussi di sviluppo software. Richiede tuttavia un'installazione e una manutenzione sistemistica non trascurabile tramite container Docker.
\end{enumerate}

La combinazione di Power BI con un applicativo di PII Discovery leggero e mirato (come l'accoppiata ARX e Microsoft Presidio) rappresenta un'alternativa concreta ed 
efficiente per mappare e proteggere le informazioni sensibili dell'azienda. Questa architettura ibrida permette di blindare la sicurezza del dato alla radice e di 
abilitare un reporting sicuro, rinunciando tuttavia ad alcune funzionalità esclusive di una piattaforma di Data Governance nativa, quali la gestione di un glossario 
aziendale centralizzato per gli utenti di business, la tracciabilità visiva del Data Lineage e gli strumenti di collaborazione diretta (apertura di task, thread di discussione 
e gestione degli incidenti qualitativi).



\section{Considerazioni sull'adozione}
A completamento del \emph{Proof of Concept} sviluppato in ambiente locale, la scelta strategica sulla futura evoluzione dello stack tecnologico di Wintech S.p.A. 
deve fondarsi su un bilancio oggettivo tra i benefici operativi e l'onere di adozione, manutenzione e gestione infrastrutturale delle diverse soluzioni analizzate.

Il primo parametro discriminante è rappresentato dai requisiti hardware e dalla complessità di installazione e manutenzione delle piattaforme di catalogazione:
\begin{itemize}
    \item \textbf{OpenMetadata}: richiede l'esecuzione simultanea di quattro container Docker (il database transazionale MySQL/PostgreSQL, il server centrale in Java, il motore di ingestion basato su Apache Airflow e il cluster di indicizzazione Elasticsearch/OpenSearch), con un consumo minimo garantito di 6-8 GB di RAM e 4 vCPU per ambienti di produzione\cite{site:openmetadata-docs, site:manuale-utente};
    \item \textbf{Open Data Discovery (ODD)}: si attesta come un candidato \emph{open-source} alternativo più leggero, richiedendo solo tre componenti (la piattaforma ODD, un database PostgreSQL e l'agente collector), con requisiti sistemistici ridotti a 4 GB di RAM e 2 vCPU. Tuttavia, non dispone di un motore nativo di Auto-Classification per i dati sensibili, vanificando parte dei vantaggi della catalogazione automatica;
    \item \textbf{Microsoft Purview}: rappresenta il modello di riferimento enterprise per l'ecosistema Microsoft aziendale, integrando nativamente l'intelligenza artificiale e la Data Compliance in \emph{cloud}, ma comporta costi di licenza e complessità di configurazione (richiedendo spesso l'uso di Kubernetes) sproporzionati rispetto alle reali necessità correnti dell'organizzazione;
    \item \textbf{Tool leggeri standalone (ARX, Presidio, Hawk-Eye)}: non richiedono un'infrastruttura server sempre attiva (24/7), potendo essere eseguiti come script Python locali schedulati periodicamente in background (tramite Windows Task Scheduler o cron in Linux), riducendo a zero l'onere sistemistico fisso e azzerando i costi di adozione.
\end{itemize}

Sulla base di queste considerazioni, l'introduzione immediata di una piattaforma di Data Catalog completa come OpenMetadata risulterebbe parzialmente sovrodimensionata per 
l'attuale contesto operativo di Wintech S.p.A., determinando un rischio di sovra-ingegnerizzazione (\emph{over-engineering}) a fronte di un patrimonio informativo che può 
essere protetto con strategie più snelle.

Tuttavia, il vero valore del progetto di stage non risiede nell'adozione immediata e vincolante di un singolo applicativo software, bensì nella maturità e nella consapevolezza 
metodologica acquisite dall'azienda in materia di Data Governance, qualità del dato e Sviluppo Sicuro. Lo studio compiuto ha evidenziato come le tematiche di conformità al 
GDPR (inclusa la corretta classificazione dei tag e l'esigenza di una tracciabilità trasparente dei flussi) non siano meri vincoli burocratici, ma requisiti architetturali 
ad alto valore aggiunto che coinvolgono trasversalmente molteplici reparti aziendali:
\begin{itemize}
    \item Il team di Ricerca e Sviluppo acquisisce strumenti concreti per bonificare e anonimizzare i database di produzione prima di utilizzarli negli ambienti locali di test e sviluppo software;
    \item Il reparto Business Intelligence ottiene linee guida chiare per progettare flussi informativi intermediati e sicuri, riducendo il rischio di esporre accidentalmente record riservati o di propagare anomalie qualitative nelle metriche direzionali;
    \item L'area Compliance e Sicurezza dispone di un \emph{framework} di riferimento (ispirato agli standard ISO/IEC 27001 e ISO/IEC 29100) per gestire con efficacia gli audit interni, mappare i rischi di re-identificazione statistica e rispondere tempestivamente alle richieste di esercizio dei diritti degli interessati (come il diritto all'oblio).
\end{itemize}

In conclusione, lo scenario evolutivo più realistico e sostenibile per l'organizzazione prevede l'adozione di un approccio modulare e progressivo: anziché mantenere attiva 
un'infrastruttura di catalogazione complessa a livello aziendale, Wintech S.p.A. potrà capitalizzare i risultati dello stage implementando in produzione gli strumenti di 
PII Discovery e anonimizzazione deterministica (come ARX) per la bonifica dei database fisici operazionali, garantendo il pieno rispetto della privacy fin dalla progettazione 
(\emph{Privacy by Design}) all'interno di flussi analitici e applicativi intrinsecamente sicuri. %8
    % capitolo 9: conclusioni
\chapter{Conclusioni}
\label{cap:conclusioni}

\intro{Il presente capitolo traccia un bilancio dell'esperienza di stage presso Wintech S.p.A., valutando i risultati ottenuti rispetto agli obiettivi 
prefissati e riflettendo sul valore apportato all'azienda e sulle competenze acquisite durante il percorso.}\\

\section{Consuntivo finale}

\section{Raggiungimento degli obiettivi}

\section{Conoscenze acquisite}

\section{Valutazione personale}
 %9

    \appendix

    \backmatter
    % \glsaddall % per stampare tutto anche se i termini non ci sono ancora nel testo

    % Crea un foglio completamente bianco (destra e sinistra) per staccare dal Cap 8 - prima di acronimi e abbreviazioni
    \cleardoublepage
    \thispagestyle{empty} \mbox{} \clearpage
    \thispagestyle{empty} \mbox{} \cleardoublepage

    \cleardoublepage
    % IMPONE IL NOME ACRONIMI o GLOSSARIO FISSO A SINISTRA NELL'INTESTAZIONE E NULLA A DESTRA
    \fancyhead[L]{\fontsize{12}{14.5} \selectfont Acronimi e abbreviazioni}
    \fancyhead[R]{} 
    \printglossary[type=\acronymtype, title=Acronimi e abbreviazioni, toctitle=Acronimi e abbreviazioni]

    \cleardoublepage
    \fancyhead[L]{\fontsize{12}{14.5} \selectfont Glossario}
    \printglossary[type=main, title=Glossario, toctitle=Glossario]

    % STESSA COSA PER BIBLIOGRAFIA E RINGRAZIAMENTI (se dovessero occupare più di una pagina)
    \cleardoublepage
    \fancyhead[L]{\fontsize{12}{14.5} \selectfont Bibliografia}
    \cleardoublepage
\chapter{Bibliografia}

\nocite{*}

% Print book bibliography
\printbibliography[heading=subbibliography,title={Riferimenti bibliografici},type=book]

% Print report bibliography
\printbibliography[heading=subbibliography,title={Report e studi},type=report]

% Print site bibliography
\printbibliography[heading=subbibliography,title={Siti web consultati},type=online]


    % --- INIZIO A SINISTRA PER I RINGRAZIAMENTI (eccezione) ---
    \cleardoublepage      % Va alla prima pagina di destra disponibile
    \thispagestyle{empty} % Rende la pagina di destra completamente bianca
    \mbox{}               % Crea la pagina invisibile
    \clearpage            % Salta alla pagina di sinistra successiva (pari)
    
    \fancyhead[L]{\fontsize{12}{14.5} \selectfont Ringraziamenti}
    \cleardoublepage
\phantomsection
\pdfbookmark{Ringraziamenti}{ringraziamenti}

\chapter*{Ringraziamenti}

Desidero innanzitutto ringraziare il mio relatore, [Nome Cognome], per la disponibilità, i preziosi consigli e la guida costante durante la stesura di questo elaborato: 
il suo supporto è stato fondamentale per portare a termine questo percorso nel migliore dei modi.
 
Un ringraziamento speciale va a Wintech S.p.A., per avermi accolta e avermi dato l'opportunità di mettere alla prova sul campo quanto imparato in questi anni di studio. 
In particolare ringrazio [Nome tutor aziendale], per la pazienza, la fiducia e il tempo dedicato a seguirmi durante tutto lo stage, rendendo questa esperienza non solo 
formativa, ma anche piacevole e stimolante.
 
Un grazie di cuore va alla mia famiglia, che mi ha sostenuta con affetto incondizionato in ogni momento di questo percorso, gioendo con me nei successi e incoraggiandomi 
nei momenti di difficoltà. Senza il vostro sostegno, tutto questo non sarebbe stato possibile.
 
Infine, ringrazio i miei amici e compagni di corso, con cui ho condiviso ansie, soddisfazioni, notti prima degli esami e infinite risate: avervi accanto ha reso questo 
cammino molto più leggero, e i ricordi di questi anni resteranno tra i più cari.
 
A tutti voi, grazie di cuore.

\bigskip
\noindent\textit{\myLocation, \myTime}
\hfill \myName

    % pagina aggiuntiva bianca alla fine
    \clearpage
    \thispagestyle{empty} % Toglie intestazione e numero di pagina
    \mbox{} % Scatola invisibile per costringere LaTeX a creare la pagina

\end{document}